%%%%% FORMAT THE PAPER %%%%%%%%%%%%%%%%%%%%%%%%%%%%%%%%%%%%%%%%%%%%%%%%%%%%%%%%

%\documentclass[draft,10pt]{thesis}
\documentclass[final,12pt]{thesis}
%\documentclass[compact,10pt]{thesis}
%\documentclass[final,10pt,sflabels]{thesis}
        % Wisconsin thesis requires a single-sided printout, and the page numbers
        % should always at the up right corner ('oneside'), see thesis.cls
        % Curly brackets: Specify the "thesis" class file (thesis.cls).
        %   It assumes the existence of a few other style (*.sty) files.  
        %   See the thesis.cls file for details.
        % Square brackets, specify the options you want to use.
        %   Allowed options are 10pt, 11pt, 12pt, draft, & final.
        %   See the "thesis.cls" file for descriptions.

% NOTE: If you want your figure captions, table captions, and bibliography
% items single-spaced (this is permissible for UW theses!), you can use
% the "\fixspacing" command in each of these environments.  This command
% will choose the appropriate spacing, whether you use the "final" or
% "draft" modes.  See the "thesis.cls" class file for details.
% CLARIFICATION: The \fixspacing command has to go inside EACH of your
%     \tablenotetext blocks to be effective

%\usepackage{lscape}     % Specify this package for landscape-style tables.
% You can use this by wrapping your table in a \begin{landscape}--
% \end{landscape} environment.  Note that the rotated table will
% probably *not* show up in Xdvi, but *will* show up in the PostScript
% version (e.g., can be seen with gv).



\usepackage{xspace}
%\usepackage{longtable}   % Specify this package for multipage tables

%\usepackage[final]{graphicx}    % Use the LaTeX2e graphics package
\usepackage{graphicx}    % Use the LaTeX2e graphics package

\usepackage{subfig}   % Specify this package to create multipart figures
% with multiple files combined in a single numbered figure 
% i.e., Fig. 1(a), Fig. 1(b), etc.

\usepackage{amsmath,amssymb}  % Load the package for extended mathematical symbol set
\usepackage{amsfonts} % Provides \mathcal
%\usepackage{axodraw}  % Necessary for using the ZEUS psfigadd command
%\usepackage{placeins}
%\usepackage{caption2}
%\usepackage{epsfig}
%\usepackage{wrapfig}
\usepackage{rotating} % Provides sidewaysfigure
%\usepackage{mcite}
%\usepackage{amscd}
%\usepackage{array}
%\usepackage{multirow}
%\usepackage{supertabular}
%\usepackage{dcolumn}
\usepackage{xspace}
%\usepackage{color}
\usepackage{calc}
\usepackage{ifthen}
%\usepackage{cite} % NEED THESE???
%\usepackage{notoccite} % too

\usepackage{tikz} % feynman diagrams!
\usetikzlibrary{trees}
\usetikzlibrary{decorations.pathmorphing}
\usetikzlibrary{decorations.markings}

\newcommand{\pb}{ pb\textsuperscript{-1} }
\newcommand{\lunits}{$cm^{-2} s^{-1}$ }
\newcommand{\dsigdM}{$ \frac{d\sigma}{dM_{inv}} $ }
\newcommand{\sieie}{$ \sigma_{i \eta i \eta } $ }
\newcommand{\detain}{$ \Delta \eta _{in} $ }
\newcommand{\dphiin}{$ \Delta \phi _{in} $ }
\newcommand{\dCotTheta}{$ \Delta \cot \theta $ }
\newcommand{\pt}{$ p_{T} $ }
\newcommand{\Et}{$ E_{T} $ }
\newcommand{\DR}{$ \Delta R $ }
\newcommand{\Zg}{$ Z/ \gamma^{*} $ }
\newcommand{\Zee}{$ Z \rightarrow e^{+} e^{-} $ }
\newcommand{\qqZee}{$ q \bar{q} \rightarrow Z \rightarrow e^{+} e^{-} $ }
\newcommand{\qqZgee}{$ q \bar{q} \rightarrow Z/ \gamma^{*} \rightarrow e^{+} e^{-} $ }
\newcommand{\Zgee}{$ Z/ \gamma^{*} \rightarrow e^{+} e^{-} $ }
\newcommand{\Ztautau}{$ Z/ \gamma^{*} \rightarrow \tau^{+} \tau^{-} $ }
\newcommand{\Wenu}{$ W \rightarrow e \nu $ }
\newcommand{\Wtaunu}{$ W \rightarrow \tau \nu $ }
\newcommand{\ttbar}{$ t \bar{t} $ }

%%%%%%%%%%%%%%%%%%%%%%%%%%%%%%%%%%%%%%%%%%%%%%%%%
%
%  Common definitions
%
%  N.B. use of \providecommand rather than \newcommand means
%       that a definition is ignored if already specified
%
%  so I changed them all...
%
%%%%%%%%%%%%%%%%%%%%%%%%%%%%%%%%%%%%%%%%%%%%%%%%%

% ---------- Physics symbols
\providecommand{\pT}{\ensuremath{p_{\mathrm{T}}}\xspace}
\providecommand{\eT}{\ensuremath{E_{\mathrm{T}}}\xspace}
\providecommand{\pt}{\ensuremath{p_{\mathrm{T}}}\xspace}
\providecommand{\et}{\ensuremath{E_{\mathrm{T}}}\xspace}
\providecommand{\kt}{\ensuremath{k_t}\xspace}

\providecommand{\ppbar}{\ensuremath{{p\overline{p}}}\xspace}
\providecommand{\bbbar}{\ensuremath{{p\overline{p}}}\xspace}
\providecommand{\ttbar}{\ensuremath{{p\overline{p}}}\xspace}


% ---------- Luminosities
\providecommand{\Lllow}  {\ensuremath{\mathcal{L}=\text{10}^\text{31}\,\text{cm}^\text{$-$2}\,\text{s}^\text{$-$1}}\xspace}
\providecommand{\Llow}  {\ensuremath{\mathcal{L}=\text{10}^\text{32}\,\text{cm}^\text{$-$2}\,\text{s}^\text{$-$1}}\xspace}
\providecommand{\Lmed}  {\ensuremath{\mathcal{L}=\text{10}^\text{33}\,\text{cm}^\text{$-$2}\,\text{s}^\text{$-$1}}\xspace}
\providecommand{\Lhigh}  {\ensuremath{\mathcal{L}=\text{10}^\text{34}\,\text{cm}^\text{-2}\,\text{s}^\text{$-$1}}\xspace}

\providecommand{\pbinv} {\mbox{\ensuremath{\,\text{pb}^\text{$-$1}}}\xspace}
\providecommand{\fbinv} {\mbox{\ensuremath{\,\text{fb}^\text{$-$1}}}\xspace}
\providecommand{\lumi}{\ensuremath{\mathcal{L}}\xspace}
\providecommand{\Lumi}{\ensuremath{\mathcal{L}}\xspace}%both upper and lower

% ---------- Units
\providecommand{\MeV}{\ensuremath{\,\text{Me\hspace{-.08em}V}}\xspace}
\providecommand{\GeV}{\ensuremath{\,\text{Ge\hspace{-.08em}V}}\xspace}
\providecommand{\TeV}{\ensuremath{\,\text{Te\hspace{-.08em}V}}\xspace}
\providecommand{\mev}{\ensuremath{\,\text{Me\hspace{-.08em}V}}\xspace}
\providecommand{\gev}{\ensuremath{\,\text{Ge\hspace{-.08em}V}}\xspace}
\providecommand{\tev}{\ensuremath{\,\text{Te\hspace{-.08em}V}}\xspace}


% Define styles for the different kind of edges in a Feynman diagram
\tikzset{
%  photon/.style={decorate, decoration={snake}, draw=npurple},
%  photon/.style={decorate, decoration={snake}, draw=purple},
  photon/.style={decorate, decoration={snake}, draw=black},
%  boson/.style={decorate, decoration={snake}, draw=npurple},
%  boson/.style={decorate, decoration={snake}, draw=purple},
  boson/.style={decorate, decoration={snake}, draw=black},
%  zboson/.style={dashed, draw=black},
  zboson/.style={decorate, decoration={dash}, draw=black},
%  electron/.style={draw=jblue, postaction={decorate},
%  electron/.style={draw=blue, postaction={decorate},
  electron/.style={draw=black, postaction={decorate},
    %        decoration={markings,mark=at position .55 with {\arrow[draw=jblue]{>}}}
  },
%  fermion/.style={draw=jblue, postaction={decorate},
%  fermion/.style={draw=blue, postaction={decorate},
  fermion/.style={draw=black, postaction={decorate},
    %        decoration={markings,mark=at position .55 with {\arrow[draw=jblue]{>}}}
  },
%  gluon/.style={decorate, draw=wred,
%  gluon/.style={decorate, draw=red,
  gluon/.style={decorate, draw=black,
    decoration={coil,amplitude=4pt, segment length=6pt}} 
}



% Line numbers for drafts, thanks to Carlos P.d.l.H. from AMANDA.
% To use this feature, say \linenumbers right after \begin{document}.
% Be sure to remove before printing out the final version!
%                                       --Katherine Rawlins, Oct. 2001
%\usepackage[displaymath,mathlines]{lineno}


%\nofiles       
        % Uncomment to prevent the creation of new *.toc *.lot, and *.lof 
        %   files.  Do this if you need to edit these files for the final
        %   version of the paper.  For example, if you have multi-page 
        %   figures or tables but you want the lists of figures & tables
        %   to have only one line for these, you will need to delete the
        %   extra lines in the *.lof and *.lot files and re-run LaTeX.

%\graphicspath{{figs/}}  % Specifies a subdirectory containing graphics files
                        % Multiple directories possible, each with inner {}
                        %   within the outer {}; e.g., {{figs/}{~/plots}}


   % Use \includeonly{} if you are just working on one \include{}'d
   % section.  This prevents LaTeX from re-processing the unneeded
   % chapters, but preserves references and page/fig/table numbers.
   % Try it.  You'll like it. 
   % Default is to compile entire document.
%\includeonly{introduction}

%\input{define}          % Read in list of user-defined LaTeX commands
%\include{citations}
%\include{zeus_def}

\begin{document}        % BEGIN THE DOCUMENT TEXT

%  The cool linenumbers feature... comment it out for the final version
%\linenumbers



%%%%% BEGIN TITLE/CONTENTS PAGES %%%%%%%%%%%%%%%%%%%%%%%%%%%%%%%%%%%%%%%%%%%%%%

\mytitle{
Z boson cross section measurement using CMS at the LHC
}
\myauthor{Jessica Lynn Leonard}
\myyear{2011}

% These aren't really needed except for the UMI abstract
\adviser{Wesley H. Smith}
\adviserrank{Professor}

\ttlpage                % This command produces the title page.  Nifty!
\cpypage                % This produces the copyright page.  Also nifty!
                        % Reminder: it costs money to retain copyright
                        
% The following three lines are for UMI abstract only
                        % Reminder: UMI abstract is required by gradschool but it is NOT a part of your thesis
%\begin{umiabstract}
%\input{umiabstract}    % *.tex file for UMI abstract page
%\end{umiabstract}      % Need to use \input rather than \include to avoid
                        %  implicit \clearpage's.  You shouldn't have any
                        %  cross-references between abstract and text, so
                        %  \input should be OK.

\frontmatter            % Choose roman-type numbers; reset page numbers to 1.

\pagestyle{thesis}      % Set the page style to "thesis" 
                        %   (page numbers at the top right)

%%%%% BEGIN ABSTRACT ETC. %%%%%%%%%%%%%%%%%%%%%%%%%%%%%%%%%%%%%%%%%%%%%%%%%%%%%
%% \begin {abstract}

%% \noindent Inclusive dijet and trijet production in deep inelastic $ep$ scattering has been
%% measured for $10<Q^2<100$ GeV$^2$ and low Bjorken $x$,
%% $10^{-4}<x_{\rm Bj}<10^{-2}$.  The data were taken at the HERA $ep$ collider
%% with center-of-mass energy $\sqrt{s} = 318$ GeV using the ZEUS detector
%% and correspond to an integrated luminosity of $82~{\rm pb}^{-1}$.  Jets were
%% identified in the hadronic center-of-mass (HCM) frame using the $k_{T}$ cluster
%% algorithm in the longitudinally invariant inclusive mode.  Measurements of dijet and
%% trijet
%% differential cross sections are presented as functions of $Q^2$, $x_{\rm Bj}$, jet
%% transverse energy, and jet pseudorapidity.  As a further examination of low-$x_{\rm Bj}$ dynamics,
%% multi-differential cross sections as functions of the jet correlations in transverse
%% momenta,
%% azimuthal angles, and pseudorapidity are also presented.  Calculations at $\mathcal{O}(\alpha_{s}^3)$ generally
%% describe the trijet data well and improve the description of the dijet data compared
%% to the
%% calculation at $\mathcal{O}(\alpha_{s}^2)$.

\noindent The cross section of 
$p p \rightarrow Z \rightarrow e^+ e^-$ 
was measured for $\sqrt{s} = 7$ TeV $pp$ collisions 
at the Large Hadron Collider using the 
Compact Muon Solenoid detector.  
Well-reconstructed electrons within the region 
$ |\eta| < 1.4442, 1.566 < |\eta| < 2.5$ 
and with a transverse energy of at least 25 GeV 
were used, 
with an invariant mass window of 60 to 120 GeV.  
In 36.1 \pb of data, 8453 candidate events were found.  
The cross section was measured to be 
990 $\pm$ 48 pb, 
which agrees well with the NNLO prediction 
of 972 pb.  

\end{abstract}
      % *.tex file for abstract
%% %\include{sponsor}      % *.tex file for "sponsorship" page
%% %\include{dedicate}     % *.tex file for dedication
%% %\include{quotes}       % *.tex file for quotations
%% \begin{acknow}
%ACKNOWLEDGEMENTS HERE

I would like to thank the professors 
who have mentored and advised me 
during my time as a grad student, 
Wesley Smith and Sridhara Dasu, 
as well as the scientists I have worked with 
over the years: 
Pam Klabbers, Sascha Savin, and Monika Grothe.  
I would also like to thank my fellow CMS students 
for their camaraderie and willingness to share knowledge, 
and my friends in Madison and in Geneva 
for providing a social outlet and a listening ear.  

Finally, I would like to thank my family, 
who never stopped encouraging me 
throughout this journey.  

\end{acknow}
   % *.tex file for acknowledgements

\tableofcontents        % *.toc file with table of contents
\listoffigures          % *.lof file with list of figures   % SWITCHED
\listoftables           % *.lot file with list of tables    % THESE TWO
     % Note that the \listoftables does not play well with the
     % deluxetable environment.  You will probably have to do some
     % hand-tweaking of the master.lot file if you have multi-page
     % tables. 
\clearpage
\pagebreak              % Enforce the start of a new page


%\prelimpages % from the UWash template -- what does it do?



%%%%% BEGIN TEXT OF PAPER %%%%%%%%%%%%%%%%%%%%%%%%%%%%%%%%%%%%%%%%%%%%%%%%%%%%%

\mainmatter

%\setfoot{\small \itshape \centerline{ DRAFT \today }} 

% This is a hand-tweak by Dolan to fix his odd landscape tables in
% compact style.  REMOVE THIS for your thesis!
%\ifkompakt\renewcommand\arraystretch{1.1}\else\fi

%\textpages  % UWash thesis template has this, what is it?

%\include{chapter_1_intro}                % *.tex file for the Introduction
%\include{chapter_2_evolution}                % *.tex file for the Introduction
%\include{chapter_3_detector}      
%\include{chapter_4_LO_MC}
%\include{chapter_5_NLO}
%\include{chapter_6_reconstruction}
%\include{chapter_7_selection}
%\include{chapter_8_analysis}   
%\include{chapter_9_results}
%\include{chapter_10_conclusion}    % *.tex file for the Conclusion

\chapter{Overview}
\label{over}
\section{Measurement of Cross Section}
\label{over:xsec}

   * particles?  although that's theory intro (but need SOME particle-specific stuff) 

%the practical purpose of MC -- can simulate all stages of gen/sim/reco/etc and compare 
%everything step-by-step with data, as well as of course getting some quantities 
%that require extra knowledge, like acceptance (and predicting efficiencies, rates, etc)
%HA, that sounds like a few of those things in the pythia manual intro!  and I came 
%up with them myself!  

%glossary? (explanatory text, with true glossary in appendix?)  

Introductory terminology?  put terminology section in beginning?

want to talk about Z and why important, motivations

   * events
In particle physics, an ``event'' is a particle interaction 
that has been captured and recorded by the detector.  
%   * signal, background.  
In general, %events that contribute to the specific process     
events caused by the specific process     
being studied are termed ``signal,'' 
while any events 
%that do not contribute 
from other sources 
are called ``background.'' 
A significant part of any analysis is designing 
criteria that select signal events while 
rejecting background events, 
so that the analysis focuses only on data containing 
that process.  
The number of events for a given signal can be used 
to calculate quantities of interest, 
for example how often that process  
occurs relative to others.  

   * cross section!
The cross section, the quantity being calculated in this analysis, 
is such a measure of ``how often things happen.''  
However, it's not measured as a rate (occurrences per unit time), 
it's measured as a cross-sectional area 
and represents the probability of the given interaction occurring.  
An analogy may be made to trying to hit a target with a tennis ball.  
The bigger the target is, the more likely you will be able to hit it.  
The sizes of both objects matter: 
a trajectory that would cause a tennis ball to just miss the target 
would cause a basketball to clip the edge, 
solely because of the basketball's larger size.  
Therefore ``cross section'' may be more precisely defined as the 
effective cross-sectional area of the target, 
taking into account the sizes of both the target and the projectile.  
When both the target and projectile are particles, %such as protons, 
they may interact ``at a distance,'' that is, without ``touching'' each other, 
through the fundamental forces.  
(For example, electrons are thought to be point particles 
and therefore have no spatial extent, 
but they still attract and repel other particles 
through the electromagnetic force.) 
This interaction-at-a-distance increases the effective cross section.  
An interaction involving two particles that 
interact with each other very strongly has a large cross section.  
The cross section depends on the strength of the force between them.  
FORMULA AND EXPLANATION HERE??

According to the definition, ``interaction cross section'' applies only to 
the particles doing the colliding.  
However, a cross section is usually associated with the entire process, 
such as (in our case) 
$ q\bar{q} \rightarrow Z/ \gamma^{*} \rightarrow e^{+} e^{-} $.  
Here, the number calculated as the cross section takes into account 
only those events where a \Zg is formed, 
as well as the fact that only some of those events decay into electrons.  
The fraction of events that decay to a certain final state is known as 
the branching ratio, abbreviated BR.  
When it is not explicitly mentioned along with the cross section 
for a given final state, it is understood to be included.  
Adding up the cross sections for all of these possible interactions 
would give the full proton-proton interaction cross section.  

%xsec(Z) = n(Zee) / ( L * BR(Zee) * eff )
\[
\sigma_{ Z } \times \textrm{ BR }_{ Zee } = \frac{ n_{ Zee } }{  \mathcal{ L } \times \epsilon \times A}
\]
where
\[
%\epsilon = \epsilon_{ acc } \times \epsilon_{ trig } \times \epsilon_{ reco } \times \epsilon_{ sel }
\epsilon = \epsilon_{ trig } \times \epsilon_{ reco } \times \epsilon_{ sel }
\]

   * acceptance, efficiency (especially for that formula above)


   * ``data-driven'' stuff?  need for efficiency, but also systs??
The term ``data-driven'' is used to describe an analysis method 
that uses only real data, taken with the detector, 
instead of relying on simulated data.  
The purpose of using a data-driven method is to 
eliminate the possibility of a physics result 
being affected by an error or inaccuracy 
in the simulation itself.  
Since some quantities, such as efficiency, 
are much easier to measure using simulated data, 
devising data-driven methods to measure those quantities 
is an important and potentially challenging 
part of analysis.  

\section{Anatomy of Event}

%Like, what actually happens!

%bunches of protons collide, maybe two protons interact very strongly
%define parton

%A detector for a particle accelerator 
%records 
A proton-proton interaction begins 
with beams of protons accelerated 
in opposite directions.  
The protons are formed into ``bunches,'' 
with many, many millions of protons per bunch.  
The beams of proton bunches are directed to 
cross at the center of the detector.  
%Most such ``bunch crossings'' do not result 
%in a significant interaction.  % interaction at all? YES, DUH pileup
Each such ``bunch crossing'' results in 
a very small number of proton interactions, 
the vast majority of which do not interact 
strongly enough to be interesting.  
However, sometimes a parton, 
i.e. a quark or gluon inside the proton, 
interacts very strongly with a parton 
from another proton.  
These are the types of interactions 
that are typically interesting.  
%   * ``hard'' vs ``soft''
%A ``hard'' interaction is one in which the partons 
%interact very energetically.  
An interaction in which the partons interact 
very energetically is a ``hard interaction.''  
Conceptually, the partons ``hit each other hard.''  
Hard interactions are typically detected 
by having end-product particles with 
a lot of momentum 
in the transverse direction, 
perpendicular to the protons' original 
direction of motion.  
Only in hard interactions is the original 
momentum disturbed so much; 
in ``soft,'' or low-energy, interactions, 
most of the protons' momentum 
continues in the same direction, 
down the beam-pipe.  

%exchange a particle
%produce another particle which then decays
Oftentimes the two partons exchange 
a particle, such as a photon or gluon.  
In other cases another particle is formed, 
such as a Z boson; 
this is the scenario 
studied in this analysis.  
Particles formed in this way 
are typically heavy and therefore 
short-lived, 
decaying into lighter, more stable particles.  
This analysis studies the Z's decay into two 
electrons, 
which are the lightest charged particles 
and therefore (to conserve charge) 
do not decay.  
These decay products are what fly ``out'' 
into the detector with 
some energy and direction, 
and it is these quantities 
that are measured by the detector.  

However, additional processes may contribute 
to the signature left in the detector.  
%
%ISR, FSR
One of the initial partons 
or final decay products 
may radiate an additional particle 
which then ends up in the detector.  
This is known as ``initial-state radiation'' (ISR) 
and ``final-state radiation'' (FSR) respectively.  
In addition, there are two ways that 
particles unrelated to the hard interaction 
may show up in the detector: 
underlying event and pileup.  
%
%   * underlying event!
The ``underlying event'' refers to the 
interactions taking place between the 
proton remnants, 
the parts of the proton ``left over'' 
from the hard interaction.  
%
%   * pileup.  maybe under, like, ``the anatomy of an event'' or ``the structure of an event?''
``Pileup'' refers to any interactions that happen 
between other protons in the bunch.  
These are generally soft, 
low-energy interactions.  
Statistically, at most only one hard interaction 
happens in any given bunch crossing.  
Both underlying event and pileup can contribute 
energy deposits that are recorded as part of the event.  
Typically they are both fairly low-energy and 
therefore easily filtered out as background.  

\section{Terms in particle physics, technical stuff}

%   * ``event''!  

   * ``high-energy''
``High-energy'' physics refers to modern particle physics. 
In the beginnings of particle physics, 
the available accelerators operated at much lower energies.  
Low-energy accelerators are still used in nuclear physics 
and aid in materials physics and biophysics studies. 
However, to probe the fundamental interactions between particles, 
today's accelerators use much higher energies.  

   * ``kinematic'' maybe previous section?
In the following pages, much will be described by the word 
``kinematic'' or ``kinematics''.  
This refers to the properties of position or motion.  
So, ``kinematic criteria'' means selections applied according to 
quantities such as momentum or angle.  


   * histogram?  errors? (stat/syst) etc?

   * statistical and systematic uncertainties or errors

   * ``phase space''! % -- move this to only place (so far) 
%that talks about it?  but theory WILL talk about it
``Phase space'' refers to the collection of total 
possible values for a set of quantities.  
A given system is represented by one set of 
values out of that collection.  
%Direction (in three dimensions)   % NOT ACTUALLY TRUE!!!
%plus speed or energy define a 
%four-dimensional phase space 
%for a particle's kinematic quantities.  
For some calculations it is necessary to include 
contributions from the entire phase space, 
for example integrating over all possible 
angles for a particle's direction.  

\section{coordinates used?}

   * center of momentum reference frame

For some high-energy physics experiments, 
there is a choice of several reference 
frames that can be used for measurement, 
each with their own benefits.  
%
However, in the case of proton-proton collisions, 
the logical reference frame for measurements is 
the frame at rest with respect to the detector, 
the ``lab frame,'' 
because it is the same as the 
proton-proton center of momentum frame.  
Therefore all measurements are done 
in the lab frame.  

%   * eta, phi, transverse quantities -- relate eta to y and all that, 
%talk about difference, talk about meaning of each one

A particular set of coordinates is often used to describe 
the direction of outgoing or intermediate particles.  
The direction of the incoming particles defines an axis, 
the beam axis.  
The azimuthal angle, measured around this axis 
and analogous to longitude on the earth's surface, 
is labeled $\phi$.  
The angle along the axis, the polar angle 
(analogous to latitude), 
is labeled by $\theta$.  
However, it is much more common to measure this 
angle in terms of pseudorapidity, $\eta$, 
which is given as a function of $\theta$:
\[
\eta = -\ln \left[ \tan \left( \frac{\theta}{2}\right) \right]
\]
$\eta$ is preferred because of its relation to another quantity, 
rapidity, $y$, which is a function of the particle's 
energy and momentum:
\[
y = \frac{1}{2} \ln \left( \frac{E+p_L}{E-p_L} \right)
\]
where $p_L$ is the longitudinal component of the 
particle's momentum (the component parallel to the beam axis).  
% WON'T DO COMPOSITION OF VECTORS
Rapidity is useful when dealing with relativistic speeds, 
such as those at which the particles typically travel.  
When viewed from separate reference frames that are moving 
at relativistic speeds with respect to each other, 
a particle's speed, direction, and energy appear to be different.  
However, rapidity has special status in that the 
difference in rapidity 
between two particles does not change 
when the reference frame is changed.  
A particle's rapidity can be approximated by the 
pseudorapidity if it is very light or 
traveling very fast 
(pseudorapidity is identical to rapidity 
if the particle has zero mass, 
in which case it is also 
moving at the speed of light).  
In addition, outgoing particles tend to be 
distributed uniformly in terms of rapidity, 
as opposed to the polar angle $\theta$, 
which does not enjoy the same uniform distribution. 
Since rapidity itself is not a measure of angle, 
pseudorapidity is used as the coordinate.  
Rapidity is calculated in cases 
where the particle of interest may be heavy 
and may not have a high speed, 
for example the Z boson examined in this analysis.  

\section{General Analysis programs}

%   * ROOOOOOOOOT!  other software, like CMSSW? yes

% root.cern.ch duh

% have no idea how to reference CMSSW -- workbook??

Two software programs in particular were 
used in this analysis.  

The CMS Software (CMSSW) package is 
the collaboration-developed software 
used for many aspects of CMS operation and use.  
It encompasses data-taking, reconstruction, 
simulation, and analysis.  
% what else to say?  

ROOT is a general data analysis package, 
written in C++, 
and developed and maintained at CERN. 
It is widely used in the field of 
experimental particle physics.  
ROOT implements many tools and functions 
necessary to various types of analysis, 
such as histograms, fitting, 
and statistics.  
CMSSW can interface with ROOT, 
providing access to ROOT's %wide 
%range of 
capabilities within a CMSSW 
analysis setting.  


\chapter{Theory Introduction}  % should probably change this name
\label{theory}

\section{The Standard Model}
\label{theory:SM}

% quarks, leptons, bosons!  HIGGS ha ha
% make separate paragraphs for different types?  or sections?  that's how other people did it

% MAKE TABLE SHOWING WHAT FEELS WHAT FORCE!  or something -- quantum number is more complicated way of showing it.  draw connection in text

% EXPLAIN SQUARE ROOT OF S so do need to get into s.  okay, via s-level diagram

The Standard Model describes the current understanding 
of how fundamental particles interact.  
There are three different types of particles: 
quarks and leptons, 
which make up matter, 
and bosons, which transmit forces.  
Quarks and leptons are collectively called fermions; 
%because they have spin-$\frac{1}{2}$   % HAVE TO EXPLAIN SPIN AND ALL THAT???
two fermions cannot exist in the same state as each other.  
Bosons do not have this constraint; 
many bosons can be in the same state.  
The particles are shown in Figure~\ref{fig:StandardModel}. 
The first three columns show the first three generations 
of the quarks and leptons, 
while the last column shows the bosons.  
Each fermion generation contains two quarks, 
a neutrino, and a charged lepton.  
In general, the generations get progressively heavier.  
The electron is lighter than the muon, 
which is lighter than the $\tau$.   
The same relation exists between the quark generations; 
the top quark is the heaviest particle in the Standard Model.  % CHECK!!!!!
Normal, everyday matter is made up solely of particles 
from the first generation.  
All normal matter is made of atoms, 
which in turn consist of a nucleus and 
electrons surrounding the nucleus.  
Atomic nuclei are made of protons and neutrons, 
which are made of up and down quarks: 
the proton by $uud$ and the neutron by $udd$.  

There are six quarks, often represented by their first letters: 
up, down, charm, strange, top, and bottom.  
The first quark in each pair has charge $+\frac{2}{3}$, 
%(in terms of the electron charge), 
while the second has charge $-\frac{1}{3}$.  
These charges are given in terms of the electron charge; 
though the charges are fractional, 
the quarks are always arranged into composite 
particles that have integer multiples of the electron charge.  

The leptons consist of the electron, muon, and tau ($\tau$), 
each with charge -1, 
and their corresponding neutrinos, 
the electron neutrino, muon neutrino, and $\tau$ neutrino, 
which are chargeless.  
Until recently it was believed that all the neutrinos are massless.  
However, recent experiments indicate 
that neutrinos do have a small mass.  % REFERENCE

The bosons in the last column transmit the forces 
by which particles interact.  
%The massless photon ($\gamma$) and gluon ($g$) 
%transmit the electromagnetic and the strong forces, 
%respectively, 
The massless photon ($\gamma$) transmits the 
electromagnetic force; this particle is more familiar in its 
everyday context as the particle that makes up light.  
The gluon ($g$), also massless, transmits the strong force, 
which holds protons and neutrons together     %%%%%% STRENGTH OF FORCES, mass of propagator?? but not in all cases
(like glue, hence the name).  
The relatively heavy $Z$ and $W$ bosons 
transmit the weak force, 
which causes nuclear decay.  

 \begin{figure}[htb]
  \begin{center}
%    \includegraphics[width=360pt]{Figures/theory-standard-model-symmetry-mag.png}
    \includegraphics{Figures/theory-standard-model-symmetry-mag.png}
  \end{center}
  \caption[Particles of the standard model]
	  {Particles of the standard model that are known to exist. 
	    The Higgs boson is also predicted as part of the 
	    Standard Model, 
	    but has not yet been discovered. 

	    }
  \label{fig:StandardModel}
 \end{figure}

%antiparticles
In general, each particle has an antiparticle partner 
of the opposite charge and quantum numbers; % opposite quantum numbers in general, explain quantum numbers? because have to account for photon, gluon, Z, neutrinos
quantum numbers are used to describe the properties of the particles 
and corresponding interactions.  
For instance, all leptons have a non-zero ``lepton number;'' 
if the initial state of a set of particles 
has a total lepton number of 1, 
then any interaction between those particles will 
result in a final state with the same lepton number.  
Lepton number is always conserved in this way, 
as are charge and other quantum numbers.  
The antiparticle of the electron ($e^-$) 
has the opposite charge and lepton number 
and is called the positron ($e^+$).  
Neutrinos, being chargeless, 
have antineutrinos with opposite lepton number.  % EXPLAIN!!!!!
Quarks and gluons have another quantum ``number'' known as color.   % EXPLAIN HOW CANNOT HAVE BARE QUARKS
The three colors are red, green, and blue, with the corresponding 
``anti-colors'' being anti-red, anti-green, and anti-blue.  
%This ``color'' does not refer to our everyday experience of red, green, and blue, 
%%but there is a certain analogy to light in the ``quantum color'' behavior.  
%but the name is not accidental -- there is an analogy to colored light.  
Color is unique in that all observed states of quarks and gluons 
are ``colorless.''
In two-quark states, a color and its anti-color ``combine'' to make a colorless state.  
Three-quark states are also colorless, but through a different combination: 
the three quarks are colored red, green, and blue 
(or the corresponding anticolors), which ``combine'' (like colored light) to form 
a state with no color.
  
Not shown in Figure~\ref{fig:StandardModel} is the Higgs boson, 
which is predicted by the Standard Model 
to give mass to the other particles.  
The Higgs has not yet been discovered; 
finding it (if it exists) is one of the 
goals of the LHC.  

The Standard Model does not include the final 
fundamental force, gravity.  
Gravity is the weakest of the forces and does not 
measurably affect interactions between fundamental particles.  
The reason that it is the only one observable on a 
cosmological scale is because its ``charge,'' mass, 
is cumulative; there is no anti-mass to counter it 
the way a negatively-charge electron and a positively-charged 
proton combine to make an electrically neutral hydrogen atom.  % put this sort of stuff in beginning?
Gravity has been described on a macroscopic scale, but 
%so far gravity has not yet been (fully?) worked into the % FIXME!!!
%framework of the fundamental particles.  
so far no theory has fully united gravity with the 
other fundamental forces.  
%There are many efforts currently ongoing to unite them, % REFERENCE?
%but none of these approaches have yet been proved.  
Such a ``theory of everything'' is 
%a dream of 
one of the goals that 
theoretical physics %.  
pursues.  

% http://www.symmetrymagazine.org/cms/?pid=1000064  the symmetry magazine SM picture
% don't know how it prints in black and white
% LATEX DOESN'T DO GIFS

TABLE OF MASSES




gauge groups, little intro to group theory

A ``group'' is a set of elements with specific properties.  
In particular, any transformation of an element in the group 
results in another element in the group.  

A ``Lie'' group is one in which any transformation % NEED THIS????
can be done in very small steps, 
each of which results in an element still in the group.  
The different fundamental forces are described by different Lie groups.  

which particles feel which forces?

Each of the three forces in the Standard Model 
is described by such a Lie group.  % group. 
%In each case, the ``dimension'' of the group 
%is the dimension of the matrices
Each group is represented by a set of square matrices 
of some dimension, given by the number in parentheses.  
The matrices used for these representations are 
all unitary (hence the common ``U'' designation).  
``Unitary'' means that the matrix's inverse 
is equal to its complex conjugate transpose.  
In other words, if you flip the matrix's elements 
about its diagonal and replace each imaginary unit $i$ 
with $-i$, then multiply it by the original matrix, 
the result is the identity matrix.   % EXPLAIN IDENTITY MATRIX??

%EM: U(1)
The electromagnetic force is described by the $U(1)$ group.  
Essentially, this describes rotation about an axis.  % talk about more than matrices?
The single parameter (an angle) requires only one matrix, 
corresponding to a single % force particle(s)
force carrier, the photon.  
The electromagnetic force causes interactions between 
photons and any particles that have electric charge.  
This includes quarks, leptons, and the charged weak bosons.  
%Because the photon is massless, the electromagnetic force % EXPLAIN MORE???
%can act over long distances, observable on the human scale.  

%Weak: SU(2)
The weak force is represented by the $SU(2)$ group.  
The ``S'' addition in the name stands for ``special'' 
and means that each matrix's trace 
(the sum of elements along the matrix's diagonal) is zero.  
$SU(2)$ represents rotations in three dimensions 
using three matrices, 
and therefore has three parameters corresponding 
to the single angle of $U(1)$.  
These give rise to the three carriers of the weak force, 
the two oppositely-charged W's and the Z.  %%% FIGURE OUT HOW TO DO W+- in latex? sort of like writing it out longhand, but should show at least W+ and W- notation 
%All three carriers have significant mass, which means 
%the distance over which they have an effect is very small.  
The weak force acts on any particles that have the ``weak 
version'' of electrical charge, 
namely all fermions.  
In particular, the weak force is the only way 
neutrinos interact.  
%However, interactions with W's only happen in ways  % TALK ABOUT CONSERVATION OF CHARGE IN GENERAL
%that conserve charge.  

%QCD: SU(3)
Finally, the strong force is represented by $SU(3)$, 
which does not have an easy analogy to rotations 
like the previous two groups.  
It uses eight matrices, which correspond to the 
eight different gluon states that mediate the force.  
% talk about dim-3 = 3 colors?
The three dimensions of the group relate to the 
three ``color charges,'' red, green, and blue.  
Strong force interactions happen between all particles 
that carry this color charge, 
namely quarks, but also gluons themselves.  
% talk about color confinement and asymptotic freedom?  YES, but how to explain in terms of mass???
This so-called self-interaction on the part of the gluons 
causes the strong force to increase with distance.  
At short distances, the quarks making up another particle 
essentially act free; 
only when they get further apart do they feel the force 
keeping them in their bound state.  
This situation is related to the fact that quarks are 
never observed alone, only within these bound states.  
An interaction that is energetic enough to cause 
quarks to separate is also energetic enough to 
create more quarks with which the original quarks 
form new bound states.  
%When a quark gains a significant energy 
%(through an interaction with another particle), 

%This is known as color confinement
%The implication is that in spite of the fact that 
%gluons are massless, 
%the strong force is not observable over large distances.  

However, the individual groups by themselves 
do not necessarily represent the reality of particles.  
It is known experimentally that the bosons carrying 
the weak force have mass.  
%The $SU(2)$ formulation of the weak force
In order for the theory to correctly predict the 
fact those particles should be massive, 
the electromagnetic and weak interactions had to be combined 
into a single framework, 
%Steven Weinberg and Abdus Salam 
%The electromagnetic and weak forces were combined into one framework, 
$SU(2) \times U(1)$, using a specific method.  
The integrated $SU(2) \times U(1)$ group 
allows for the collective four force-carrying bosons 
(photon, positively- and negatively-charged W's, and Z) 
to be described as compositions of the same underlying states.  
The fact that the three weak bosons are massive, though, 
comes from applying ``spontaneous symmetry breaking.''  
This method uses the fact that the groups mentioned 
above have underlying symmetries: 
that is, they contain parameters that do not affect the 
physical description in any way.  
They are symmetric with respect to variations in these parameters.  
However, the values of these parameters can be chosen 
in a strategic way and the result manipulated % TALK ABOUT FIELD THEORY AT ALL?? this is all lagrangians and mass terms
to give masses to the weak bosons.  
(A consequence of this ``strategic choice'' is the appearance 
of another massive particle in the theory: 
the Higgs boson.  
The Higgs boson is therefore predicted by the theoretical 
mechanism used to account for massive weak bosons.  
Time -- and collision data -- will tell if the theory 
in this form is correct.)  


%Transformations within these groups can be represented as matrices.  % IS THIS ACTUALLY THE SAME AS THE GROUP MATRICES??
These matrices act on the set of initial states 
and result in the set of final states.  
Each element in the matrix, mapping a single initial state 
to a single other final state, 
gives the amplitude with which the % ONLY IN SCATTERING???  make sure this is right context 
transformation happens.  
Essentially, squaring the matrix element gives the 
relative probability of that particular transformation.  
Matrices are often used to represent transitions 
from one quantum mechanical state to another.  

This is used in particle physics as such: 
a transformation from an initial state, 
in our case two quarks from colliding protons, 
to a final state, 
an electron and a positron, 
is represented by a Feynman diagram, 
shown in Figure~\ref{fig:ZeeFeynmanDiagram}.  
The interaction can be thought to progress to the right.  
Each straight line represents a fermion, 
either the quark or the electron.  
The wavy line represents a boson, the Z.  
(Gluons are represented by curly lines.)  
The intersections between particle lines, or vertices, 
show interactions between the given particles.  
%The type of interaction can be determined by the particles   % NEED THIS?  
%taking part: 
%an interaction between a Z and a quark or lepton 
%indicates the weak force is involved.  
The matrix element corresponding to this transformation 
can be written down 
by attributing factors to the various elements 
of the Feynman diagram, 
the particle lines and the vertices.  

somewhere get into spinors...

put equation for differential cross section element 
in terms of amplitude/matrix element 



put all qcd stuff in here?  don't need all the detail of zeus stuff

\section{Electroweak Physics}
\label{theory:EWK}

lagrangian etc? i.e. more specifics on ewk physics

feynman diagrams, how they can actually be used -> 
matrix elements -> cross section element.  

\begin{figure}[htb]
\begin{center}
  \begin{tikzpicture}[
      thick,
      % Set the overall layout of the tree
%      level/.style={level distance=1.4cm, line width=0.8mm},
      level/.style={level distance=1.4cm, line width=0.5mm},
      level 2/.style={sibling distance=1.4cm},
      level 3/.style={sibling distance=1.4cm}
    ]
    \coordinate
%    child[grow=south east]{
%      child[grow=north east]{
%        edge from parent [gluon]
%        node[right]{$g$}
%      }
    child[grow=south east]{
      edge from parent [electron]
      child[grow=south west] {
        edge from parent [electron]
%        node[above] {$\bar{q}$} % ADDED
        node[below] {$\bar{q}$} % ADDED
%          child[grow=south west]{
%            edge from parent [electron]
%            node[above] {$\bar{q}$}
%          }
%          child[grow=south east]{
%            edge from parent [gluon]
%            node[right] {$g$}
%          }
      }
      child[grow=east, level distance=2.4cm] {
        child[]{ % added the [], did nothing, don't know how to make angle of line the same as LHS
          edge from parent [electron]
          node[below] {$e^{-}$}
        }
        child[]{ % added the [], did nothing, don't know how to make angle of line the same as LHS
          edge from parent [electron]
          node[above] {$e^{+}$}
        }
%        edge from parent [boson]
        edge from parent [zboson]
%        node[below] {$Z/\gamma*$}
        node[below] {$Z$}
      }
%    }
      edge from parent [electron] node [above=3pt] {$q$}
    };
  \end{tikzpicture}
\end{center}
%  \caption[Feynman diagram of \qqZgee]
  \caption[Feynman diagram of \qqZee]
%	  {Feynman diagram of \qqZgee.  
	  {Feynman diagram of \qqZee.  
	    }
  \label{fig:ZeeFeynmanDiagram}
\end{figure}

where to put cross section formula stuff 
-- make new section?
(and which stuff to use?)

order of diagram = number of vertices!  

\subsection{Z Production}
\label{theory:Zprod}

s, inv mass, pdfs attach to tree-level matrix element, 
Z/gamma* interference

s = Z mass -- well, dunno, also includes momentum

%The definition of a 
A particle system's ``invariant mass'' is 
a combination of its energies and momenta that 
remain unchanged regardless of reference frame.  
A particle (or everyday object) can appear to 
have a different energy depending on the reference 
frame from which it is viewed.  
If you are sitting on an airplace at cruising altitude, 
the plane is not moving with respect to your body.  
However, if you are on the ground, the plane is 
moving very fast -- 
it has a different kinetic energy in that frame.  
The physics of the situation cannot depend on 
the reference frame of the observer, 
even at the relativistic speeds of interacting particles.  
The invariant mass is such a quantity that 
remains the same.  
In addition, since the energy and momentum 
of a system are conserved, 
the invariant mass is also conserved 
throughout the entire interaction, 
in particular the decay.  
The general definition is given as 
\[
M_{inv}^2 = \left( \sum E \right)^2 - \left\| \sum \mathbf{p} \right\|^2
\]
For a single particle, the expression simplifies 
to the particle's mass.  
This analysis makes use of the fact that the 
invariant mass of the decay particles 
(the two electrons) 
is the same as the (invariant) mass of the 
particle which decayed, the Z.  
The above formula, simplified to two particles, 
\[
M_{inv} = \sqrt{ \left(E_1 + E_2\right)^2 - \left\|\mathbf{p}_1 + \mathbf{p}_2\right\|^2 }
\]
reconstructs the mass of the Z boson 
based on the energies and momenta 
of the decay products.  

Another diagram contributes to the process being studied. 
In this case, a photon replaces the Z boson 
as the link between the quarks and the electrons.  
Because these diagrams have the exact same 
initial and final states, 
they ``interfere'' with each other.  
Experimentally, it is impossible to tell which one happened; 
both must be included in the calculations.  % is this exactly why???

It is necessary to note that the photon that takes 
part in this interaction is a ``virtual'' photon -- 
%it cannot be observed by itself.  
it cannot be directly observed.  
The reason for this is that there is always a 
center-of-momentum reference 
frame for the two particles colliding to form the photon.  
In other words, there is always some way that you, 
the observer, can be moving, 
such that it looks like the two particles 
have the same speed but opposite directions.  
In this reference frame, the photon would be formed at rest, 
which violates the principle that the speed of light 
is the same in all reference frames -- 
the speed of light can never be zero.  
Equivalently, such a photon would have mass, 
while real photons are massless.  
The mass of this virtual photon is indeed the same 
as the invariant mass of the system.  
%Such virtual particles can exist due to inherent   %% is that actually why these photons in particular can be virtual?  
%fluctuations in energy, by the uncertainty 
%principle of quantum mechanics.  
Such virtual particles are allowed to exist 
by the uncertainty principle of quantum mechanics, 
which governs the ``borrowing'' of energy from the vacuum.  
The smaller the time scale over which the energy 
fluctuation happens, 
the larger the energy fluctuation allowed.  

Having a specific mass defines a thin 
surface or ``shell'' in kinematic phase space; 
only certain combinations of energy and momentum 
are allowed, 
because of the relations between mass, energy, and momentum.  
Particles that are virtual, 
i.e. have a mass much different than their 
defined ``real'' mass, 
are called ``off-shell'' -- 
they are ``off the mass shell.''  

The Z boson, while it has a defined mass, 
%DECAY WIDTH ETC ETC ETC 
can be ``off-shell'' in the same way as the photon.  
That is, in general the 
%reconstructed 
mass 
does not take on a single precise value; 
instead, over many events it manifests as a distribution 
peaked around the ``real'' value.  
This is typical of particles that are measured only 
by their decays; 
these particles tend to be at least slightly virtual 
%because they are not being observed directly.  
because of their short lifetimes.  
%in fact, 
%In addition, the width of the peak is inversely related to the   % how relevant is this actually?  
%lifetime of the particle.  

PIC OF DY SPECTRUM??

In this case, only the area around the Z peak is being 
studied, not the full spectrum.  
The main contribution in this area is the Z itself 
as opposed to virtual photons.  


\subsubsection{Results from Previous Experiments}
\label{theory:prev}

\subsection{Z Decay}
\label{theory:Zdec}
need this section??


somewhere have intro to QFT and all that?  

also, explanatory items from overview:

   * [here or] in theory chapter define ``tree-level''.  
ALSO IN THEORY: explain how Feynman diagrams are so useful! 
writing down the matrix element and all that
AND PDFs, basically whole process of stuff simulated 
really happens and should be explained.  
AND DEF OF PARTONS
yeah, basically explain all the stuff in the MC chapter
INCLUDING ISR and FSR
ALSO define ``jets'' and talk about how they come from 
quarks and gluons (define partons) and relate to ``hadronization'' -- 
not entirely relevant for this analysis, 
but possible in general
AND define ``color''
AND asymptotic freedom, how it takes more and more energy to get them further apart

   * intial- and final-state radiation

   * [HERE OR IN THEORY INTRO] need to do on-shell/off-shell decays 
to explain why Z mass has a spectrum and not a single value

   * invariant mass (here?)

   * cross section formula and matrix elements and what said matrices do

   * PDFs, what they are, how they relate to Z physics

AND, of course read the others' theses to see what they had in here, too.  




\chapter{Experimental Setup}
\section{CERN}
International laboratory in Europe, founded 1959?, maybe little significant history.  

\section{LHC}
All the good stats on the LHC: circumference, (place,) dipole design, magnets, magnetic field.  Design energy and luminosity.  

\section{CMS}
Important stats on CMS: height, weight, etc.  

Figure: Event display, Fig. \ref{fig:EventDisplay}

 \begin{figure}[htb]
  \begin{center}
    \includegraphics[width=360pt]{CMS-BW.pdf}
  \end{center}
  \caption[Three-dimensional display of a \Zee event in the CMS detector]{Three-dimensional display of a \Zee event in the CMS detector.}
  \label{fig:EventDisplay}
 \end{figure}

\subsection{Detection of Particle Interactions}
\subsection{Coordinate System}
\subsection{Magnet}
\subsection{Luminosity}
Figure: Luminosity plot, Fig. \ref{fig:LuminosityVsTime}

 \begin{figure}[htb]
  \begin{center}
    \includegraphics[width=360pt]{CMS-BW.pdf}
  \end{center}
  \caption[Luminosity as a function of time]{Luminosity as a function of time.}
  \label{fig:LuminosityVsTime}
 \end{figure}

\subsection{Calorimeters}
\subsubsection{ECAL}
\subsubsection{HCAL}
\subsection{Muon Systems}
\subsection{Tracker}
\subsection{Pixels}
\subsection{Trigger}
\subsubsection{Level-1 Trigger}
General L1 information, specific electron algorithms in online selection section.  
\subsubsection{High-Level Trigger}
General HLT information, specific electron algorithms in online selection section.  

\chapter{Event Simulation}
\label{sim}
%somewhere have the pretty picture from 
%matt mc tutorial on how generators/geant/etc 
%fit in with everything else (reco, whatever)

%ARXIV PAPER OF AWESOMENESS
%http://arxiv.org/abs/hep-ph/0403045

%PYTHIA MANUAL OF AWESOMENESS
% the one I'm looking at is lutp0613man2.pdf, 
% but there's probably a more recent one




% I THINK I HAVE TOO MANY LONG SENTENCES 
% AND PREPOSITIONAL PHRASES


%need for event simulation: 

Within %the role of % CHANGE THIS CONSTRUCTION
a high-energy physics experiment's role of making 
discoveries and measurements, 
it necessary to know ahead of time what exactly 
should be expected. 
What will the detector signature of a new process look like? 
How will we know whether we've seen something unexpected? 
This is where event simulation contributes.  
%what exactly is the simulated data
Essentially, a series of programs is used to carefully 
generate and calculate all the relevant quantities for a set of 
fake events.  
This information is then used 
to aid analysis in ways for which 
real data alone cannot suffice.  

Having the ability to simulate physics processes 
serves multiple purposes. 
It aids in detector design: 
knowing the expected typical characteristics of particle 
interactions is essential to design a detector 
suited for those interactions.  
Once the detector design is settled, 
simulation of the detector is useful to design 
the algorithms used to reconstruct particles 
from their signature interactions with the 
detector material.  
In addition, simulating the physics processes 
can give estimates of how many 
events of a particular type are expected, 
further aiding the design process.  
The event simulation also contributes 
directly to many analyses, 
in the way of calculating the acceptance, 
the fraction of events that can theoretically 
be detected (see Section~\ref{FIXME}): 
it is impossible to know from observation 
how many events are missed by the detector, 
because of the very fact that they are unseen.  
It is instead necessary to get that fraction 
from a framework in which 
the characteristics of all events are 
inherently known, not reconstructed.  
Finally, % and possibly most visibly (?) I don't like that
simulated data is directly compared 
with ``real'' data from the detector 
to interpret the real-data results.  
If the data shows something significantly 
different from the simulation, 
then something is missing: 
perhaps a calibration needs to be applied, 
or the response of a detector unit needs 
to be further understood.  
Or, perhaps, there is a new physics process 
appearing for the first time, 
which was not previously known and 
which was therefore not present in the simulation.  
In this case simulation of proposed new physics 
processes may narrow down the identity of the 
observed new process.  
Whatever the case, 
discrepancy between observed data and the event 
simulation indicates that further 
investigation is necessary.  


%points of event simulation: 

%   * directly relevant: compare with data (description of current understanding as is) 
%-- we want to be able to interpret what we're seeing 
%in terms of what physics process it might be

%   * also relevant: acceptance is necessarily a MC calculation!

%   * also: design detectors, design reconstruction and other (e.g. trigger) algorithms

%   * also: event rates (like trigger rates?)  

The entire detection process is simulated, 
including the protons' direct interaction 
and any subsequent particle decays, 
as well as how the end-product particles 
interact with the detector as they pass through 
and how the detector itself responds.  
This latter part includes not only the 
material of the individual subdetectors 
but also the algorithms of 
the Level-1 Trigger, 
which are implemented in hardware 
(see Section~\ref{exp:L1}).  
The High-Level Trigger does not need to be 
simulated in this way; 
since its algorithms are all software-based, 
the same code can be run without modification 
on both real data and simulated data.  



\section{Monte Carlo Event Generation}
\label{sim:MC}
put why called ``monte carlo'', 
also connect ``event simulation'' terminology 
with ``generator'' and the multiple types thereof.
like, ``simulation'' broken down into several steps, 
one of which is the generation of the main 
physics process itself (aka ``hard physics process'')
HOW the events are actually ``generated'', 
like the whole random number business. 
It uses a random number generator , 
to determine ``in a probabilistic way'' -- 
only one main interaction happens 
for each proton (bunch) interaction, 
but the sum over all events ,
according to the probability of each type of 
interaction happening.  



\subsection{Put this stuff in its own section?}
\label{sim:MCexplain}

%types of programs with different types of output -- what they're used for. 
%start with tree-level stuff and dsigma eqn...

%   * cross section integrators

%   * event generators

There are two sorts of programs that can be called 
%``Monte Carlo generators''.  
``Monte Carlo event simulators''.  
They both use the matrix elements derived from 
the Feynman diagrams of the physics interactions of interest 
(explained in Section~\ref{FIXME}), % THAT MEANS EXPLAIN ALL THIS STUFF IN PREVIOUS CHAPTER
but they use the matrix element formalism 
in different ways to achieve different ends.  
As a reminder, the equation to obtain the 
differential cross section % WHAT'S IT ACTUALLY CALLED??
from the matrix element for a given diagram is 

\[
d\sigma = \frac{1}{\hat{s}}|\mathcal{M}|^2 \frac{d \cos \theta d\phi}{8(2\pi)^2}
\]

The first type of Monte Carlo program outputs 
a single quantity or set of quantities characteristic 
of the given process, 
such as the cross section.  
These programs are called ``cross section integrators''.  
They calculate 
$d\sigma$ for each of a uniform distribution of 
phase space vales (i.e. uniform increments 
of $\cos\theta$ and $\phi$).  
$d\sigma$ represents each event's ``weight''; 
the total cross section $\sigma$ is then obtained 
by integrating over the full collection of $d\sigma$ values.  
The final result is a value for the cross section; 
hence the name.  
However, since the events themselves are distributed 
uniformly in $\cos\theta$ and $\phi$, 
their characteristics 
do not represent actual event distributions 
(real event distributions in data have 
non-uniform distributions, 
in $\cos\theta$ in particular).  
Therefore they can only be used to calculate 
a few select quantities; 
they can't be used to predict distributions.  

The second type of program are called ``event generators''.  
They also use the matrix element to calculate 
the differential cross section % REAL NAME
with the equation given above.  
However, instead of using $d\sigma$ as an event weight, 
it ``unweights'' the events to obtain the physical 
distributions.  
It does this through the so-called acceptance-rejection method: 
first the maximum possible value of $d\sigma$ is calculated, 
$d\sigma_{max}$. 
Each event's individual $d\sigma$ value, 
the ratio $\frac{d\sigma}{d\sigma_{max}}$ is taken; 
this number is necessarily a fraction between 0 and 1.  
A random number $g$ between 0 and 1 is then generated.  
If the ratio $\frac{d\sigma}{d\sigma_{max}}$ is greater 
than $g$, then then event is kept.  
Otherwise, it is rejected.  
This method causes the generated events to have the properties 
and distributions that events would have in reality -- 
sections of phase space with a higher ``event weight'' 
are more likely to be represented.  

%tree-level all well and good, but need corrections to really 
%make things realistic.  
%ways of doing the corrections:

%   * matrix element

%   * parton shower

However, using a single Feynman diagram is not typically 
sufficient to make very accurate calculations.  
%a number or determine distributions is not typically 
%sufficient.  
In general, there are significant contributions from 
the higher-order diagrams  
(diagrams with more legs and loops, 
but still representing the same final arrangement of particles). 
These corrections must be taken into account.  
In addition, the matrix elements and Feynman diagrams 
for these calculations only consider bare quarks, 
that is, quarks extracted from the framework of 
hadrons such as the proton.  
These bare quarks cannot exist in reality, 
so a method of accounting for the quarks 
as being inside hadrons must also be used.  
There have been two different paths taken to deal % path = tack? track?...
with the above issues.  
In the first, the ``matrix element'' method, 
Higher-order diagrams are explicity included in 
the matrix element calculation.  
Since the higher the order, 
the more complex the calculation, 
only diagrams of up to a given order are 
generally used.  
In one implementation, more ``legs'' can be dealt with, 
but diagrams containing virtual loops are ignored.  % EXPLAIN IN PREVIOUS CH
In the other, all diagrams, including those with virtual 
loops, of a single order are used, 
but at the moment there exist satisfactory solutions 
only for next-to-leading-order diagrams 
(or NLO, explained in Section~\ref{FIXME}).  
%only for diagrams one step higher 
%than the basic tree-level diagram
The matrix element method does not explicitly 
take care of the bare-quark problem; %, though.  
this can be handled in multiple ways.  
%Programs who do not deal with the interactions of 
Some cross section integrators deal with the interactions of % is it only cross section integrators??
quarks inside hadrons using Feyman's factorization theorem. 
Programs which do not deal with them in this way 
can be interfaced with programs following the 
second of the two paths, %.  
%The second path is called 
the ``parton shower'' method, 
which is the method used by many event generators 
(as opposed to cross section integrators).  
In the parton shower method, 
instead of including diagrams representing 
higher orders and quarks-within-hadrons 
into the matrix element calculation, 
this part of the interaction is taken care of 
by a separate routine.  
This routine takes the products of the 
basic tree-level interaction 
and decays them according to 
probabilities and rules coded into the program.  
The process is iterative, 
repeating for any quarks or gluons that 
branch off at any point.  
In this way the parton shower method 
is more flexible than the matrix element method -- 
it can deal with many branchings and legs. 
However, it is also less precise formally, 
since it avoids the use of the matrix elements 
from the Feynman diagrams.  

%where do the PDFs fit in?  just to parton shower method?? 
%no, it looks like PDFs have to enter in both 
%ways of doing hadronization. 
%(actually, is it only just the initial 
%proton interaction where you need the PDFs?? 
%inside a proton is the only place 
%where x is a relevant quantity...
%dunno, skirt for now)

%do I need more detail at this point?  
%I guess still need to explain modular structure... 
%can more specific stuff like hadronization 
%model go just under pythia, or should it 
%go here?  
%I guess should be more general, 
%since those algorithms are more general.  
%SO, HERE

The philosophy of parton-shower generators 
in particular 
is that the entire process of simulating 
a physics event can be done in a modular fashion.  
Fig.~\ref{fig:MCsteps} 
shows the typical steps necessary to simulate an events.  
The incoming protons are modeled in terms of the 
parton distribution functions (Section~\ref{FIXME}).  % NEED TO EXPLAIN ``PARTONS'' IN THAT SECTION
The ``hard subprocess'' refers to the main 
physics interaction.  
Not shown in the diagram are other interactions 
involving different parts of the protons 
or other protons in the same bunch; 
these interactions are not as interesting 
because they are typically not as energetic.  
The participant particles in the main interaction 
are then ``evolved'' through the showering process, 
which includes emissions of photons or gluons 
by the initial or final particles, 
as well as ``branchings'' of final particles. % INCLUDED IN FORMER?? 
Any resulting quarks and gluons are 
formed into ``colorless'' hadrons, 
which may then decay into other particles. 
These final particles are the ones seen 
in the detector.  
The fact that the full process can be separated 
into these steps 
makes it easier for different programs 
to handle individual steps.  
For example, one program may simulate the 
hard process, 
while another is better-equipped to take care 
of the showering and hadronization, 
and a third is specialized in the decay of 
a particular particle.  
To further the theme of modularity, 
the fully-decayed situation 
is then plugged into a program simulating 
the passage of the particles through the detector 
and the detector's response (Section~\ref{FIXME}).  

%And how these two things fulfill the different needs for 
%both higher orders and hadronization

%BIG PRETTY PARTON SHOWER PICTURE but also need to explain PDF part, etc. etc. 
should compare to other one from zeus people, which one is better? 

 \begin{figure}[htb]
  \begin{center}
    \includegraphics{Figures/mc-partonShower-lesHouches.png}
  \end{center}
  \caption[Diagram of typical steps used to simulate event in parton shower method]
	  {Diagram of typical steps used to simulate event in parton shower method.
	    The incoming protons are modeled in terms of the 
	    parton distribution functions. 
	    The ``hard subprocess'' refers to the main 
	    physics interaction.  
	    The participant particles in the main interaction 
	    are then ``evolved'' through the showering process, 
	    which includes emissions of photons or gluons 
	    by the initial or final particles, 
	    as well as ``branchings'' of final particles. % INCLUDED IN FORMER?? 
	    Any resulting quarks and gluons are 
	    formed into ``colorless'' hadrons, 
	    which may then decay into other particles. 
	    These final particles are the ones seen 
	    in the detector.  
	  }
  \label{fig:MCsteps}
 \end{figure}


and lund string model pictures (both, yay!)

 \begin{figure}[htb]
  \begin{center}
%    \includegraphics{Figures/mc-partonShower-lesHouches.png}
    \subfloat[]{\label{fig:MClundA}\includegraphics{Figures/mc-lundStringModel-zeusPeople.png}}
    \subfloat[]{\label{fig:MClundB}\includegraphics{Figures/mc-lundStringDiagram-zeusPeople.png}}
  \end{center}
  \caption[Schematics showing the Lund string model of hadronization]
	  {Schematics showing the Lund string model of hadronization.}
  \label{fig:MClund}
 \end{figure}



pythia is definitely parton shower, 
explain steps and hadronization model, etc

%\subsection{some explanation of the black box?}
%\label{sim:MCBlackBox}
% it's not really a black box anymore!
%can talk about black-box mentality in introduction
%modular structure here?  that pretty picture showing all the 
%different steps to an interaction that need to be 
%simulated. 
%all those steps can be done by separate packages, 
%which is the philosophy.  (pythia manual)
%pdfs ->
%hard scatter (underlying event, multiple interactions in parallel) ->
%[parton] shower (both isr and fsr, incl QED) ->
%hadronization (not directly applicable to my process, but applicable for bg's) -> 
%decay resonances -> interactions with detector material ...

\subsection{Monte Carlo Generator Programs}
\label{sim:MCGens}

short intro on programs: pythia, powheg, tauola, fewz

HOW THE PROGRAMS USED TO GENERATE THESE SAMPLES AND NUMBERS 
FIT INTO THIS FRAMEWORK.  
   pythia, powheg, tauola, fewz... else??  (didn't end up using MCatNLO) 
I don't think anything else.  
also explain about fewz different (xsec integrator, sounds like)

I feel like I remember seeing something from herwig?? in one of the 
production config fragments...

\subsubsection{PYTHIA}
\label{sim:MCGensPythia}

PYTHIA THE MONSTER

what are the different tunes for?  tune z2, tune d6t

\subsubsection{POWHEG}
\label{sim:MCGensPowheg}
why is this one used in favor of pythia?  
more to the point, what exactly is it? 
and how is it different from pythia? 

\subsubsection{Other Generators}
\label{sim:MCGensOther}
Like, whichever other samples I check (background etc) -- actually, they're all pythia

TAUOLA AND FEWZ

AND, FEWZ webpage was on frank petriello's wisconsin space and doesn't exist anymore

http://www.phys.hawaii.edu/~kirill/FEHiP.htm links to 
http://www.hep.wisc.edu/~frankjp/FEWZ.html (bad link)

FEWZ 2.0 abstract on the arXiv: http://arxiv.org/abs/1011.3540 
from ryan gavin, ye li, frank petriello, and seth quackenbush (15 Nov 2010)

WHICH VERSION DID VBTF USE??  look at syst AN.  
e-mail frank about documentation page??

https://indico.cern.ch/getFile.py/access?contribId=38\&resId=3\&materialId=slides\&confId=71330  tutorial from ryan giving original arXiv references (non-2.0)




\section{Detector Simulation}
\label{sim:Detector}
%\subsection{GEANT Detector Model/Modeling of Particle Interactions in Detector Material}

\subsection{GEANT Detector Model}
\label{sim:DetectorGeant}

GEANT the giant

\subsection{Level-1 Trigger Emulator}
\label{sim:DetectorL1Emul}
talk about how emulator necessary for l1, but not for hlt (duh)

\chapter{Event Reconstruction}
\label{evReco}
%a little on daq, event building here?

%Only dealing with parts of event relevant to \Zee analysis.  

The information from the detector for any given event 
gets read out over millions of channels 
from the many different subdetectors.  
These signals must be combined to provide 
meaningful physics information 
%that can be analyzed
about the interaction that took place 
within the detector, 
namely the measurable quantities 
(like direction and energy) 
of the outgoing decay products.  

\section{Detector Object Reconstruction}
\label{evReco:detReco}
%This is for the basic objects only. 
%BUT, basic object reconstruction is so integrated with 
%physics object reconstruction, maybe combine them?

The first step of the particle reconstruction 
involves translating the information from each 
detector channel, 
such as position and electronics signal, 
into useful quantities such as position in 
$\eta$-$\phi$ space and energy or momentum.  
This is done by algorithms 
mapping each individual 
detector cell to its $\eta$-$\phi$ placement 
and calculating its energy, 
which necessarily takes into account 
the position of each cell.  
The second step involves 
connecting information from different channels 
to build up a more complex picture of 
the particles that passed through the detector.  

\subsection{Electromagnetic Calorimeter Reconstruction}
\label{evReco:ECAL}
%\subsection{Electromagnetic Calorimeter Clustering}
%reconstruction of position of energy deposits, 
%rechits and that stuff: digitization, calibration?
% where to go for this stuff?? well, reco principles or
% ecal local reco to start...

The information coming directly from the ECAL 
is read out in terms of each cell's electronics response.  
The ECAL reconstruction algorithms translate this into 
an energy value, which is then calibrated 
according to the cell's energy %as well as its 
and 
$\eta$-position.  
The algorithms also calculate quantities 
relevant to the quality of the deposit, 
in particular whether or not the deposit happened 
at a time consistent with the LHC 
bunch collisions, i.e. if it was ``in time'' or ``out of time''.  
%bunch collisions, i.e. if it was ``in time''.  
This information is especially useful in 
identifying apparent energy deposits that have 
been found to come from the 
%electronics themselves; 
signal-conversion electronics theselves (the transducers); % Wesley
these fake deposits were discovered during 
commissioning and are known as ``ECAL spikes''.  
Since the spikes can imitate energy deposits from 
actual 
%decay products, 
particles, % Wesley
they must be removed from the information used 
to reconstruct electrons and other particles.  
They can be distinguished from real particle deposits 
due to their timing 
as well as the fact that the surrounding crystals 
have no trace of a deposit.  
Therefore, if a high-energy deposit has no or 
very little energy in its neighboring crystals, 
or if its timing is not consistent with a bunch 
crossing, 
it is considered a spike and taken out of the 
reconstruction.  

%\subsubsection{ECAL Spike Removal}
%ALSOOOOO need to talk about spikes and spike removal
%(since that's now automatic, but also wasn't previously. 
%then talk about in event selection, because that was 
%a criterion...? no, because it's not a criterion for 
%e.g. newest datasets)
% HAVEN'T FOUND REFERENCE FOR WHAT'S INCLUDED IN RECO
% except for, like, swiss cross and out of time, is that enough?

%\subsection{Pixel Hits} % need?

\subsection{Track Reconstruction}
\label{evReco:track}
A track is a sequence of hits in the tracker, 
reconstructed to trace the trajectory of a passing 
particle.  
The Kalman Filter algorithm is used by default for CMS tracks, 
but is not good at catching abrupt changes in the 
electron trajectory due to bremsstrahlung.  
Therefore, the Gaussian Sum Filter algorithm 
is used for electrons, 
because it is able to find 
those hits in unexpected directions.  

\subsubsection{Kalman Filter}
\label{evReco:KF}
% CMS Note 2006/041
% CMS Note 2006/026 for seeding
The Kalman Filter (KF) tracking method \cite{CMS-NOTE-2006-041} 
is seeded \cite{CMS-NOTE-2006-026}
by pairs of hits in the pixel detector.  
The seeding algorithm first searches for a hit 
in one of the outer layers of the pixel detector.  
If a hit is found, another hit is sought in a 
layer closer to the interaction point, 
looking in an $\eta$-$\phi$ window around 
the original hit.  
Only two out of the three pixel layers are 
required to have hits 
in order to maintain a high efficiency.  
Once a track seed is found, 
further hits are sought in each successive tracking layer. 
The Kalman filter method is used to predict possible 
locations for hits in the next layer based on 
the currently-known track parameters. % and their 
%uncertainties.  
Hits in that area are then sought and included 
in the track, combining each hit with the 
predicted location for that layer and weighting 
each value depending on the value's uncertainty.  
If there are multiple candidate hits in the next layer, 
a possible trajectory is calculated for each hit, 
along with a trajectory accounting for the case 
in which the hit in that layer was lost.  
This procedure is repeated for each layer: 
each of the calculated trajectories for a given layer 
is matched to hits in the next layer, 
for which new trajectories are calculated.  
To prevent the number of necessary calculations 
from growing too large, 
the number of possible trajectories at each step 
is capped at five.  
%(BUT HOW DETERMINED WHICH ONES???)
%at each layer
In the case of multiple possible trajectories 
for a given fully-reconstructed track, 
defined by two trajectories sharing more than 
half of each of their hits,
the track candidate with fewer hits 
(or, in the case of the same number of hits, 
the track with the worse fit)
is discarded.  

\subsubsection{Gaussian Sum Filter}
\label{evReco:GSF}
% CMS Note 2005/001

The Kalman Filter track-fitting does not accurately 
account for the behavior of particles in a material, 
specifically energy lost in the material 
and changes in direction due to radiation.  
The KF description is in particular inadequate 
for electrons, which 
readily interact with matter in these ways.  
%are prone to these interactions with matter.  
The Gaussian Sum Filter (GSF) algorithm 
\cite{CMS-NOTE-2005-001} was 
therefore developed:
essentially a Kalman Filter 
with a more sophisticated energy loss modeling.   
Instead of approximating the energy loss 
(which is non-Gaussian)
with a Gaussian distribution, 
which KF implicitly does, 
the GSF method models it as a 
weighted sum of Gaussian distributions.  
This greatly improves the performance 
of the track-fitting for electrons.  
%PICTURE? trk-gsf-QoverP.png shows charge-to-momentum
% ratio for KF and GSF -- but figure out 
% WHAT Q/P ACTUALLY TELLS YOU

\section{Electron Reconstruction}
\label{evReco:elec}
% https://twiki.cern.ch/twiki/bin/view/CMS/ElectronRecoPrinciples
% https://twiki.cern.ch/twiki/bin/view/CMS/SWGuideEcalRecoClustering
%here's where we start talking about electrons
%CMS Note 2006/040
As an electron 
%from an interaction 
passes through the CMS detector, 
it leaves a signature pattern in the subdetectors.  
Characteristically, electrons produce a track 
in the tracker system, 
as well as an energy deposit in the 
electromagnetic calorimeter.  
Therefore an electron is reconstructed 
by matching tracks and energy clusters 
that appear to have been made by the same particle 
\cite{CMS-NOTE-2006-040}.  
There are currently two approaches used to 
initiate electron reconstruction: 
the ``ECAL-driven'' method and 
the ``tracker-driven'' method, 
named according to which subdetector's 
information initiates the reconstruction.  
%one starts with ECAL information 
%while the other begins with information from 
%the tracker.  
Each method produces a collection of 
basic electron candidate objects, 
each consisting of a calorimeter cluster 
and hits in the first few layers of the tracker.  
These collections are then combined into one 
and used as seeds to reconstruct the electron tracks, 
producing fully-reconstructed electrons.  
%how many electrons are one vs other vs both?

\subsection{ECAL-driven Electron Seeding}
\label{evReco:ecalDrv}

The ECAL-driven electron reconstruction method 
starts with an energy deposit (or ``supercluster'') 
in the ECAL 
and searches for hits in the 
pixel system that matches the 
position of the calorimeter deposit.  

\subsubsection{Superclusters}
\label{evReco:SC}
%Hybrid and multi5x5 (hybrid doesn't start with basic clusters, 
%so probably combine detector object and physics object sections)

The calorimeter energy deposit used in the electron 
reconstruction is called a supercluster: 
a cluster of clusters.  
A single electron or photon deposits most of its energy 
in a 5x5 square of ECAL crystals, 
as shown in Fig. \ref{fig:BasicCluster}.  
However, an electron passing through the tracker material 
in a strong magnetic field produces 
bremsstrahlung and electrons from converted photons, 
causing the original electron's energy to 
have spread out in $\phi$ once it reaches the calorimeter.  
The process of creating a supercluster gathers these 
energy deposits to form the initial energy of the electron.  

 \begin{figure}[htb]
  \begin{center}
%    \includegraphics[width=360pt]{Figures/elec-BasicCluster-opaque.png}
    \includegraphics[width=180pt]{Figures/elec-BasicCluster-opaque.png}
  \end{center}
  \caption[\fixspacing Typical energy distribution in a basic cluster]
	  {\fixspacing Typical energy distribution in a basic cluster.}
  \label{fig:BasicCluster}
 \end{figure}

% need pretty picture!

The electron reconstruction currently uses two different 
superclustering methods, 
the ``hybrid'' supercluster method in the barrel, 
and the ``multi5x5'' clustering method in the endcap.  
The hybrid algorithm forms superclusters 
and basic clusters from the same energy deposit 
information, 
while the multi5x5 algorithm first forms 
a single collection of basic clusters, 
which are then separately combined into superclusters.  

%Hybrid
%``hybridSuperClusters''
%hybrid: (dynamic phi road turned OFF)
%http://cmssw.cvs.cern.ch/cgi-bin/cmssw.cgi/CMSSW/RecoEcal/EgammaClusterProducers/python/hybridSuperClusters_cfi.py?revision=1.10&view=markup
%http://cmssw.cvs.cern.ch/cgi-bin/cmssw.cgi/CMSSW/RecoEcal/EgammaClusterProducers/src/HybridClusterProducer.cc?revision=1.32&view=markup
%http://cmssw.cvs.cern.ch/cgi-bin/cmssw.cgi/CMSSW/RecoEcal/EgammaClusterAlgos/src/HybridClusterAlgo.cc?revision=1.61&view=markup


The hybrid algorithm, 
shown in Figure~\ref{fig:HybridSuperCluster}, 
starts with a list of all 
ECAL crystals with an energy greater than a threshold; 
these crystals are considered seeds.  
Beginning with the highest seed crystal, 
the area around each seed is searched for 
energy deposits in steps of one crystal in $\phi$.
For each $\phi$ step, a ``domino'' of crystals is 
formed in $\eta$: 
one crystal to each side of the central crystal, 
or, if the energy sum of those three crystals is 
greater than a threshold, 
two crystals to each side of the central crystal, 
for a total of five crystals per domino.  
The process is done for a given number of $\phi$ 
steps in both the positive and negative directions 
(currently 17 steps).  
As other seeds are encountered in this process, 
they are removed from the list of available seeds.  
Each domino with a higher energy than its two neighbors 
is considered as the seed for a basic cluster, 
and the collection of dominos formed by the search 
process is the basis for the supercluster.  
The basic clusters are formed by summing the 
energy of each local-maximum domino with 
those of its neighboring lower-energy dominos, 
excluding dominos that have already been used.  
The supercluster is then formed by summing the energies 
of the constituent basic clusters, 
with its seed considered to be the highest-energy 
basic cluster 
and its position given by the energy-weighted positions 
of the constituent clusters.  

 \begin{figure}[htb]
  \begin{center}
    \includegraphics[width=360pt]{Figures/elec-SC-hybrid-algo.png}
  \end{center}
  \caption[\fixspacing Illustration of the ``Hybrid'' supercluster algorithm]
	  {\fixspacing Illustration of the ``Hybrid'' supercluster algorithm. 
	    Dominoes with energy above a given threshold are constructed 
	    along the $\phi$ direction in both directions out from the seed crystal.  
	  }
  \label{fig:HybridSuperCluster}
 \end{figure}



%Multi5x5:
%``multi5x5SuperClustersWithPreshower''

%https://twiki.cern.ch/twiki/bin/view/CMS/ECALDPGClusterization  also talks about supercluster energy corrections
%http://cmssw.cvs.cern.ch/cgi-bin/cmssw.cgi/CMSSW/RecoEcal/EgammaClusterProducers/python/multi5x5SuperClusters_cfi.py?revision=1.2&view=markup
%http://cmssw.cvs.cern.ch/cgi-bin/cmssw.cgi/CMSSW/RecoEcal/EgammaClusterProducers/src/Multi5x5ClusterProducer.cc?revision=1.5&view=markup
%http://cmssw.cvs.cern.ch/cgi-bin/cmssw.cgi/CMSSW/RecoEcal/EgammaClusterAlgos/src/Multi5x5ClusterAlgo.cc?revision=1.9&view=markup
%http://cmssw.cvs.cern.ch/cgi-bin/cmssw.cgi/CMSSW/RecoEcal/EgammaClusterProducers/src/Multi5x5SuperClusterProducer.cc?revision=1.2&view=markup
%http://cmssw.cvs.cern.ch/cgi-bin/cmssw.cgi/CMSSW/RecoEcal/EgammaClusterAlgos/src/Multi5x5BremRecoveryClusterAlgo.cc?revision=1.8&view=markup SCORE

The multi5x5 algorithm, 
shown in Figure~\ref{fig:Multi5x5SuperCluster}, 
first creates a collection 
of basic clusters from seeds, 
which are ECAL crystals with an energy higher than 
a given threshold.  
If a seed represents a local maximum, 
then its energy deposit is combined with those 
from the crystals in the surrounding 5x5 square 
to form a basic cluster, 
ignoring any crystals that have already been 
used to form the central 3x3 part of another cluster.  
The basic clusters are then used as seeds for the 
superclustering algorithm.  
For each seed cluster with energy higher than a threshold, 
the energies of any other clusters within an $\eta$-$\phi$ 
window of the seed (longer in $\phi$ than in $\eta$)
are combined with its own energy to create the supercluster, 
leaving out clusters that have already been used to 
form another supercluster. 
The position of the supercluster is given by the 
energy-weighted position of the individual clusters.  

 \begin{figure}[htb]
  \begin{center}
    \includegraphics[width=360pt]{Figures/elec-SC-multi5x5-algo.png}
  \end{center}
  \caption[\fixspacing Illustration of the ``Multi5x5'' supercluster algorithm]
	  {\fixspacing Illustration of the ``Multi5x5'' supercluster algorithm.
	    Clusters seeded by seed crystals are grouped into 
	    superclusters.  
	  }
  \label{fig:Multi5x5SuperCluster}
 \end{figure}



%also preshower cluster info in endcap... how?
%what kind of energy corrections done? 
%or is that for energy scale section?
%at least SC corrections are in dpg clusterization twiki

\subsubsection{Pixel Hit-matching}
\label{evReco:pixMatch}

To reconstruct the full electron, 
the supercluster must be matched with a track.  
% where exactly to split ecal-driven and tracker-driven?
This is done by extrapolating the probable position of 
tracker hits from the supercluster's position, 
%given the supercluster energy.  
as illustrated in Fig. \ref{fig:PixMatch}.  
Given the supercluster energy and the fact that a 
non-radiating electron would have hit the calorimeter 
at the energy-weighted center of the supercluster's 
energy spread, 
the electron's initial (pre-radiation) trajectory 
through the tracker can be calculated 
for both the positive and negative electron 
charge cases.  
The reconstruction algorithm then looks for a matching 
hit in the first pixel detector layer 
within a $\phi$-$z$ window of this position 
for each charge case.  
If a hit is found, another hit is sought in the second 
pixel layer within a different $\phi$-$z$ window 
of the hit in the first layer.  
%This set of pixel hits is the seed(??)

%PICTURE from my old talk or something!  

 \begin{figure}[htb]
  \begin{center}
%    \includegraphics[width=360pt]{Figures/elec-pixmatch.png}
    \includegraphics[width=240pt]{Figures/elec-pixmatch.png}
  \end{center}
  \caption[\fixspacing Schematic of pixel-hit search algorithm in electron-seeding]
	  {\fixspacing Schematic of pixel-hit search algorithm in electron-seeding. 
	    For both the positive and negative electron charge cases, 
	    a possible trajectory is extrapolated back from the calorimeter deposit 
	    given the deposit's energy. 
	    A hit is sought along either trajectory in either the inner or 
	    middle pixel layer within a window ($\Delta\phi_1$). 
	    If a hit is found, a second hit is sought along that trajectory 
	    in the next pixel layer(s) ($\Delta\phi_2$).}
  \label{fig:PixMatch}
 \end{figure}

\subsection{Tracker-driven Electron Seeding}
\label{evReco:trkDrv}
% AN 2008/032 michele pioppi
In the tracker-driven electron seeding method, 
electron-finding is seeded by tracks, 
which are followed out to the 
ECAL to search for calorimeter clusters 
associated with each track.  
This method was developed to better deal with 
low-\pT electrons, for which the energy deposits 
from bremsstrahlung may be too widely separated 
to be matched,  
and non-isolated electrons within jets, 
whose energy deposits are not necessarily 
distinguishable from those of surrounding hadrons.  
For these electrons the calorimeter cluster-based 
seeding does not perform as well as for 
isolated, high-\pT electrons.  

The tracker-driven electron seeding method 
takes as input the default collection 
of tracks, obtained by an iterative usage 
of the KF algorithm.  
The method first tries to follow each track 
all the way to the calorimeter and search 
for nearby energy clusters.  
All clusters with an acceptable E/p ratio 
with respect to the track are tested, 
and the cluster closest in $\eta$-$\phi$ 
space to the track is the one chosen. 
 
Out of the tracks that were not 
well associated to any energy cluster, 
a preselection is applied to select 
those tracks that may have resulted from 
radiating electrons, 
which would not have been found 
by the cluster-matching method.  
If the track was not matched because it had 
too few tracker hits and could not be traced 
to the calorimeter, 
or if the KF algorithm did not fit the 
track well enough because of possible 
bremsstrahlung, the tracks are passed 
on to the next step of the selection.  
A reduced GSF track-fitting algorithm 
is then run on these tracks, 
in an attempt to recover the 
bremsstrahlung-induced changes in track direction. 
This reduced version of the GSF algorithm runs 
much faster than that with the default parameters, 
but the performance is slightly reduced.  

Finally, a multi-variate analysis (MVA) using 
boosted decision trees 
is performed on these candidates.  
The MVA algorithm is 
%trained 
calibrated 
on a sample of electrons from 
$b \bar{b}$ events, which are primarily non-isolated, 
and a sample of \Zee events, 
which provides high-\pT electrons.  
%List of MVA variables?: ...
The output of the MVA is a number indicating 
the likelihood of the object being an electron.   
Those objects whose MVA output is greater than a threshold 
are considered electron candidates. 
These candidates are then combined with those 
previously identified from the cluster-matching 
to produce the final collection of 
tracker-driven electron seeds.  

\subsection{Electron Track Reconstruction}
\label{evReco:elecTrk}

The full collection of electron seeds, 
both ECAL- and tracker-driven, 
are used as the basis for reconstructing 
the electron track and hence 
completing the electron reconstruction.  
As mentioned previously, 
the electron reconstruction uses the 
GSF algorithm to reconstruct tracks.  
The electron kinematics are 
assigned to be a weighted combination 
of those of the constituent supercluster 
and track, 
and the final electron collection 
is saved for further analysis.  

%%%%%%%\subsection{Z Reconstruction}

%\subsection{Electron Energy Scale and Resolution} %???

%\subsection{Electron Isolation Variables} %???? here, or in electron selection????
%\subsection{Electron Identification Variables} %???? again, here????

\section{Beam Spot Reconstruction}
\label{evReco:BS}
%maybe include just to talk about primary vertex correction
% https://twiki.cern.ch/twiki/bin/view/CMSPublic/SWGuideFindingBeamSpot
% cms note 2007/021
The region where the beams overlap and protons may interact, 
the ``luminous region'', is called the beam spot.  
Typically this area is a few centimeters long, 
with a width of 20-50 microns.  % Wesley
The beam spot location can be used as an estimate 
for the initial position of tracks measured in the detector, 
and is therefore useful in track-seeding.  
Its position %and width %how? necessary?
%in all three dimensions %how is it done for z?
can be determined from measurements of tracks in the 
general track collection 
\cite{CMS-NOTE-2007-021}.  
%This is done by applying a fit to the distance of 
%the tracks' closest approach to the nominal beam spot (0,0,0), 
%plotted as a function
This is done by making a distribution of each track's distance of 
closest approach to the nominal beam spot (0,0,0) 
as a function of the track's $\phi$ angle at its 
closest approach.  
If the beam spot is far from (0,0,0), 
the distribution looks sinusoidal.  
This distribution is then fit with a sine curve 
parametrized in the beam spot's actual 
$x$- and $y$-positions, 
and these values are extracted.  
Initially, one beam spot position was calculated    % Wesley (I think this is what he wanted)
for a full run.  
The calculation switched to using smaller run segments 
called ``luminosity sections'' (LS) 
in order to account for changes in the beam spot position 
during a run.  

\section{Primary Vertex Reconstruction} % ? probably
\label{evReco:PV}
% few twikis with some info (most don't have much):
% https://twiki.cern.ch/twiki/bin/view/CMSPublic/WorkBookVertexReco
% https://twiki.cern.ch/twiki/bin/view/CMSPublic/SWGuideVertexReco
% https://twiki.cern.ch/twiki/bin/view/CMSPublic/WorkBookOfflinePrimaryVertexFinding
% http://twiki.cern.ch/twiki/bin/view/CMSPublic/SWGuideOfflinePrimaryVertexProduction has like the most info
% possibly cms note 2006/026 section 7 (reference in kf paper) using pixels for pv reco
% also cms note 2006/029, section 3.1 - vertex finding

Since protons may collide anywhere within the region of 
beam overlap, 
the main interaction in a given event may not necessarily 
happen at the exact center of the detector.  
In particular, there may be a significant displacement 
in the $z$-direction.  
Determining the location of the main interaction, 
also called the primary vertex, 
is an integral part of reconstructing the event.  
For example, an electron may leave a deposit in 
a given calorimeter cell that has a fixed $\eta$ position 
relative to the nominal interaction point, 
the very center of the detector (0,0,0).  
However, if the interaction producing the electron 
actually takes place at a different point, 
the electron's actual $\eta$ position will be 
different relative to the interaction.  
Therefore the position of the primary vertex 
should be taken into account when calculating $\eta$.  

% GENERAL PROCESS FOR VERTICES, BUT PV SIMPLE CASE!(?)
%The process of reconstructing the primary vertex consists 
%of two steps: vertex-finding and vertex-fitting.  
%In the first step primary vertex candidates 
%are created from the tracks in the event, 
%matching each track to the interaction 
%that likely produced it.  
%In the second, 
%the resulting vertices are fit and examined 
%to determine various parameters, 
%like the goodness of the fit to the vertex.  
%BLAH, fix if actually need
%ACTUALLY, PV reco does some fitting anyhow, 
%looks like it's already in ``pv-finding'' algo
%as described in that cms note

The process of the reconstructing the primary vertex 
\cite{CMS-NOTE-2006-026}, 
\cite{CMS-NOTE-2006-029} 
involves creating primary vertex candidates
from the tracks in the event,
matching each track to the interaction
that likely produced it.
This is done by grouping the tracks according 
to their $z$-position at the point at which they pass 
most closely to the beam spot.  
The tracks are required to pass relatively close 
to the beamline, 
and they are also required to have a \pt greater 
than some threshold (with a default of 1.5 GeV).  
Tracks within a small distance (default 1 mm) of each other 
at the beam spot are considered to have come from 
the same vertex, 
and the primary vertex candidates are formed from these 
groups of tracks.  
At this point a fit is applied to each primary vertex candidate, 
and tracks found to be incompatible with the candidate are discarded 
until all tracks meet a given compatibility level.  
Finally, vertices that are incompatible with the position 
of the beam spot are discarded.  
(A second algorithm uses the beam spot as constraint on the 
position of the vertex fit itself, 
but this algorithm was not fully commissioned at the time 
of this study.)  
If no reconstructed primary vertex candidates are found in 
an event, 
a ``fake'' primary vertex is added at the position 
of the beam spot.  
Since more than one interaction per event is possible 
but generally only one is of interest (the ``primary'' interaction), 
multiple candidates are ordered according to transverse energy of 
their constituent tracks, 
and the highest of these is taken to be the primary vertex.  

\clearpage

\chapter{Event Selection}
\section{Online}
\subsection{Trigger Eras}
%\subsection{Electron Trigger}
%\subsubsection{Electron Algorithms}
%\subsubsection{Trigger Performance}
\subsection{Level-1 Trigger}
\subsubsection{Description of Algorithm}
Description of algorithm.  
\subsubsection{Performance}
Distributions, resolutions, efficiencies of Level-1 electron/photon objects.  

Figures: Trigger Electron/Photon Pt, eta, phi spectra for L1, Fig. \ref{fig:L1TriggerObjectSpectra}

 \begin{figure}[htb]
  \begin{center}
    \includegraphics[width=360pt]{CMS-BW.pdf}
  \end{center}
  \caption[Trigger Electron/Photon Pt, eta, phi spectra for L1]{Trigger Electron/Photon Pt, eta, phi spectra for L1.}
  \label{fig:L1TriggerObjectSpectra}
 \end{figure}


Figures: \DR for L1 trigger object matching, Fig. \ref{fig:L1TriggerObjectDeltaR}

 \begin{figure}[htb]
  \begin{center}
    \includegraphics[width=360pt]{CMS-BW.pdf}
  \end{center}
  \caption[\DR for L1 trigger object matching to offline]{\DR for L1 trigger object matching to offline-reconstructed objects.}
  \label{fig:L1TriggerObjectDeltaR}
 \end{figure}

Figures: L1 trigger Pt resolution plots using offline reference, Fig. \ref{fig:L1TriggerObjectResolutions}

 \begin{figure}[htb]
  \begin{center}
    \includegraphics[width=360pt]{CMS-BW.pdf}
  \end{center}
  \caption[L1 trigger Pt resolution plots using offline reference]{L1 trigger Pt resolution plots using offline reference.}
  \label{fig:L1TriggerObjectResolutions}
 \end{figure}

Figures: L1 trigger electron/photon efficiency turn on using offline reconstructed electrons and MC, Fig. \ref{fig:L1TriggerObjectEfficiencies}

 \begin{figure}[htb]
  \begin{center}
    \includegraphics[width=360pt]{CMS-BW.pdf}
  \end{center}
  \caption[L1 trigger electron/photon efficiency turn on using offline reconstructed electrons and MC]{L1 trigger electron/photon efficiency turn on using offline reconstructed electrons and MC.}
  \label{fig:L1TriggerObjectEfficiencies}
 \end{figure}


\subsection{High-Level Trigger}
\subsubsection{Description of Algorithms}
Description of algorithms.
\subsubsection{Performance}
Distributions, resolutions, efficiencies of HLT electron/photon objects.

Figures: Trigger Electron/Photon Pt, eta, phi spectra for HLT, Fig. \ref{fig:HLTriggerObjectSpectra}

 \begin{figure}[htb]
  \begin{center}
    \includegraphics[width=360pt]{CMS-BW.pdf}
  \end{center}
  \caption[Trigger Electron/Photon Pt, eta, phi spectra for HLT]{Trigger Electron/Photon Pt, eta, phi spectra for HLT.}
  \label{fig:HLTriggerObjectSpectra}
 \end{figure}


Figures: \DR for HLT trigger object matching, Fig. \ref{fig:HLTriggerObjectDeltaR}

 \begin{figure}[htb]
  \begin{center}
    \includegraphics[width=360pt]{CMS-BW.pdf}
  \end{center}
  \caption[\DR for HLT trigger object matching to offline]{\DR for HLT trigger object matching to offline-reconstructed objects.}
  \label{fig:HLTriggerObjectDeltaR}
 \end{figure}

Figures: HLT trigger Pt resolution plots using offline reference, Fig. \ref{fig:HLTriggerObjectResolutions}

 \begin{figure}[htb]
  \begin{center}
    \includegraphics[width=360pt]{CMS-BW.pdf}
  \end{center}
  \caption[HLT trigger Pt resolution plots using offline reference]{HLT trigger Pt resolution plots using offline reference.}
  \label{fig:HLTriggerObjectResolutions}
 \end{figure}

Figures: HLT trigger electron/photon efficiency turn on using offline reconstructed electrons and MC, Fig. \ref{fig:HLTriggerObjectEfficiencies}

 \begin{figure}[htb]
  \begin{center}
    \includegraphics[width=360pt]{CMS-BW.pdf}
  \end{center}
  \caption[HLT trigger electron/photon efficiency turn on using offline reconstructed electrons and MC]{HLT trigger electron/photon efficiency turn on using offline reconstructed electrons and MC.}
  \label{fig:HLTriggerObjectEfficiencies}
 \end{figure}




\section{Offline}
\subsection{Electron Selection}
Starts out with full set of objects reconstructed as electrons.  
NOT necessarily real electrons, but energy deposits consistent with those from electrons at a low-level.  
(talk about def of GSF electron etc so know what kinds of cuts have already been applied)
Further conditions are applied to select the electron objects most likely to be real electrons from a Z.  
First electrons are matched to trigger objects passing an egamma trigger with a Level-1 threshold of 5 GeV.  
A 20-GeV cut on the supercluster transverse energy is applied.  
Next comes a set of cuts designed to reject electrons from photon conversions.  
Then come cuts to distinguish real electrons from other energy deposits which may mimic electrons, 
and finally are cuts to make sure the electron is isolated, i.e. not part of a jet (which would not be coming from a Z).  

\subsubsection{Matching Offline-Reconstructed Objects to Trigger Objects}
Figures: \DR for L1 and HLT trigger object matching, Fig. \ref{fig:TriggerObjectSelectionDeltaR}

 \begin{figure}[htb]
  \begin{center}
    \includegraphics[width=360pt]{CMS-BW.pdf}
  \end{center}
  \caption[\DR for L1 and HLT trigger object matching to offline for electron selection]{\DR for L1 and HLT trigger object matching to offline-reconstructed objects.}
  \label{fig:TriggerObjectSelectionDeltaR}
 \end{figure}


\subsubsection{Conversion Rejection}
The conversion rejection step of the electron selection aims to eliminate photons that have converted into electron pairs.  
Several variables have been chosen to distinguish between conversions and prompt electrons, focusing on the track properties.  
Since prompt electrons are expected to hit the first layer of the tracker, 
any electron-associated track that is missing initial hits can be rejected as a conversion that occurred within the tracker.  
The looser requirements of the 95 percent efficient working point allow an electron to have one missing hit;
tighter working points require zero missing hits.  
Fig. \ref{fig:MissHitsConvRejVars} shows the number of missing hits for each electron.  
In addition, the conversion rejection algorithm searches for candidate electron-pair partner tracks, 
i.e. track pairs that appear to come from the same point.  
Every electron track is compared to all opposite-charge tracks within a dR cone of 0.3, and two values are calculated: 
a measure of the difference in the angles in the r-z plane between the two tracks (dCotTheta, Fig. \ref{fig:DCotConvRejVars}), 
and the distance of closest approach between the two tracks in the x-y plane (dist, Fig. \ref{fig:DistConvRejVars}).  
Electron pairs from conversions are expected to have low values for both quantities, 
since both electron tracks begin at the point of conversion and initially have the same direction as the converting photon.  
The WP95 requirements apply a threshold at 0.02 for both quantities.  
If an electron track corresponds to a track pair having values below that threshold for both variables, 
it is assumed to be a conversion and is rejected.  

%Figures: Conversion rejection variables before respective cuts, Fig. \ref{fig:ConvRejVars}

 \begin{figure}[htb]
  \begin{center}
%    \includegraphics[width=180pt, angle=90]{missHitsSeq-14sep10.pdf}
    \includegraphics[width=180pt, angle=90]{missHitsSeq.pdf}
  \end{center}
  \caption[Missing hits conversion rejection variables after previous cuts]{Missing hits conversion rejection variables after previous cuts.}
  \label{fig:MissHitsConvRejVars}
 \end{figure}

 \begin{figure}[htb]
  \begin{center}
%    \includegraphics[width=180pt, angle=90]{dCotThetaSeq-14sep10.pdf}
    \includegraphics[width=180pt, angle=90]{dCotThetaSeq.pdf}
  \end{center}
  \caption[DCotTheta conversion rejection variables after previous cuts]{DCotTheta conversion rejection variables after previous cuts.}
  \label{fig:DCotConvRejVars}
 \end{figure}

 \begin{figure}[htb]
  \begin{center}
%    \includegraphics[width=180pt, angle=90]{distSeq-14sep10.pdf}
    \includegraphics[width=180pt, angle=90]{distSeq.pdf}
  \end{center}
%  \caption[Conversion rejection variables before respective cuts]{Conversion rejection variables before respective cuts.}
  \caption[Dist conversion rejection variables after previous cuts]{Dist conversion rejection variables after previous cuts.}
  \label{fig:DistConvRejVars}
 \end{figure}


\subsubsection{Electron Identification}
Several variables to differentiate genuine electrons from other particles that may mimic electron behavior.  

Figures: Electron ID variables before respective cuts, Fig. \ref{fig:ElecIdVars}

 \begin{figure}[htb]
  \begin{center}
    \includegraphics[width=360pt]{CMS-BW.pdf}
  \end{center}
  \caption[Electron ID variables before respective cuts]{Electron ID variables before respective cuts.}
  \label{fig:ElecIdVars}
 \end{figure}



Figures: Reconstructed electron/photon Pt, eta, phi spectra after ID, Fig. \ref{fig:RecoSpectraAfterEid}

 \begin{figure}[htb]
  \begin{center}
    \includegraphics[width=360pt]{CMS-BW.pdf}
  \end{center}
  \caption[Reconstructed electron/photon Pt, eta, phi spectra after ID]{Reconstructed electron/photon Pt, eta, phi spectra after ID.}
  \label{fig:RecoSpectraAfterEid}
 \end{figure}


\subsubsection{Electron Isolation}
Conditions to require that electrons found are isolation, i.e. not contained in a jet with other particles.  
We want (``prompt''?) electrons that come directly from the Z, not from some other interaction with lots of other end products packaged along with the electron.  

Figures: Electron isolation variables before respective cuts, Fig. \ref{fig:ElecIsoVars}

 \begin{figure}[htb]
  \begin{center}
    \includegraphics[width=360pt]{CMS-BW.pdf}
  \end{center}
  \caption[Electron isolation variables before respective cuts]{Electron isolation variables before respective cuts.}
  \label{fig:ElecIsoVars}
 \end{figure}



Figures: Reconstructed electron/photon Pt, eta, phi spectra after ID + isolation, Fig. \ref{fig:RecoSpectraAfterEidEiso}

 \begin{figure}[htb]
  \begin{center}
    \includegraphics[width=360pt]{CMS-BW.pdf}
  \end{center}
  \caption[Reconstructed electron/photon Pt, eta, phi spectra after ID + isolation]{Reconstructed electron/photon Pt, eta, phi spectra after ID + isolation.}
  \label{fig:RecoSpectraAfterEidEiso}
 \end{figure}


\subsubsection{Electron Selection Performance} %wesley comment
Figures: Reco PT resolution plots using MC reference for Z MC, Fig. \ref{fig:RecoPtResolution}

 \begin{figure}[htb]
  \begin{center}
    \includegraphics[width=360pt]{CMS-BW.pdf}
  \end{center}
  \caption[Reco PT resolution plots using MC reference for Z MC]{Reco PT resolution plots using MC reference for Z MC.}
  \label{fig:RecoPtResolution}
 \end{figure}


\subsection{\Zee Event Selection}
\subsubsection{Acceptance for \Zee Events}
\subsubsection{Efficiency of Selection for \Zee Events}
Tag and Probe:  The tag and probe method is used to determine the efficiency of a given set of electron selection steps.  
In this method, a set of data events is identified which appear to be good Z->ee candidate events.  
One of the electrons in each of these events must pass 
all selection cuts and essentially be identified as a golden electron; 
%some stringent
this is the tag.  
%The other electron must 
%pass some less stringent criteria
%be identified at least as a supercluster; 
The criteria applied to the other electron are looser: it is not required to pass the selection cuts being tested.  
This is the probe electron.  
The event's status as a Z candidate is determined by enforcing a mass window cut on the invariant mass of the tag-probe system 
(this is sufficient(enough?) to select events in which both legs are real electrons -- REFERENCE??).  
The selection efficiency is then given by the fraction of probe electrons that pass the requirement being tested.  

Figures: Efficiency table and supporting plots

\subsubsection{Distribution of Z Kinematic Variables}

\chapter{Analysis Method}
\section{Comparison of Monte Carlo and Data}
Figures: Electron-positron invariant mass for trigger objects/ID only/ID+isolation, Fig. \ref{fig:InvMassAfterEachStep}

 \begin{figure}[htb]
  \begin{center}
    \includegraphics[width=360pt]{CMS-BW.pdf}
  \end{center}
  \caption[Electron-positron invariant mass for trigger objects/ID only/ID+isolation]{Electron-positron invariant mass for trigger objects/ID only/ID+isolation.}
  \label{fig:InvMassAfterEachStep}
 \end{figure}

Figures: Same-sign dielectron invariant mass for electrons passing all selection steps, Fig. \ref{fig:SameSignInvMass}

 \begin{figure}[htb]
  \begin{center}
    \includegraphics[width=360pt]{CMS-BW.pdf}
  \end{center}
  \caption[Same-sign dielectron invariant mass for electrons passing all selection steps]{Same-sign dielectron invariant mass for electrons passing all selection steps.}
  \label{fig:SameSignInvMass}
 \end{figure}


\section{Background Subtraction}
Figures: Background subtracted yield versus eta with a MZ mass window - comparison to MC, Fig. \ref{fig:YieldVsZEtaComparedToMC}

 \begin{figure}[htb]
  \begin{center}
    \includegraphics[width=360pt]{CMS-BW.pdf}
  \end{center}
  \caption[Background subtracted yield versus eta with a MZ mass window - comparison to MC]{Background subtracted yield versus eta with a MZ mass window - comparison to MC.}
  \label{fig:YieldVsZEtaComparedToMC}
 \end{figure}

Figures: Background subtracted Z PT distribution with a MZ mass window - comparison to MC, Fig. \ref{fig:YieldVsZPtComparedToMC}

 \begin{figure}[htb]
  \begin{center}
    \includegraphics[width=360pt]{CMS-BW.pdf}
  \end{center}
  \caption[Background subtracted Z PT distribution with a MZ mass window - comparison to MC]{Background subtracted Z PT distribution with a MZ mass window - comparison to MC.}
  \label{fig:YieldVsZPtComparedToMC}
 \end{figure}


\section{Estimation of Systematic Uncertainties}

\chapter{Results}
\section{Cross Section Measurement}
Cross section calculation 
\section{Comparison to Theory}
\section{Comparison to Previous Experiments}

\chapter{Summary and Conclusions}
\label{summ}

This analysis measured the cross section of the 
$p\bar{p} \rightarrow Z \rightarrow e^+ e^-$ 
interaction within an invariant mass window of 60 to 120 GeV, 
using 36.1 \pb of $\sqrt{s} = 7$-TeV $pp$ collision data.  
Each event was required to have passed the trigger 
and to have two well-reconstructed electrons, 
which resulted in 8453 events.  
The event selection efficiency was determined using the 
tag-and-probe method to be 0.610, 
and the acceptance was calculated using a 
POWHEG (NLO) Monte Carlo \Zee sample to be 0.387.  
Background was estimated from PYTHIA Monte Carlo samples, 
from a template method using signal- and background-like 
samples, 
and from a functional fit to the 
candidate invariant mass distribution.  
Systematic uncertainties were estimated 
by varying the electron energy scale, 
varying the Monte Carlo sample used for the efficiency ratio, 
and taking errors on the fit where fits were performed.  
The error on the acceptance due to theoretical uncertainties 
was taken from the official CMS analysis.  

The cross section was determined to be 
\[
\sigma(p \bar{p} \rightarrow Z/\gamma *) \times \mathrm{BR}(Z \rightarrow e^+ e^- )
= 990 \pm 40 \mathrm{(lumi)} \pm 11 \mathrm{(stat)} \pm 17 \mathrm{(theory)} \pm 17 \mathrm{(syst)} \mathrm{pb}
\]
This result is in agreement with the FEWZ (NNLO) 
prediction of 972 $\pm$ 40 pb.  
In addition, it agrees with the CMS official measurement of 992 pb, 
as well as with the ATLAS measurement of 972 pb.  

The agreement of these results bodes well for CMS: 
the detector is well calibrated and can produce 
a solid measurement of standard model physics.  
This points toward a future of successful new 
measurements and exciting physics.  

% And I have learned a HECK of a lot.  Mission accomplished, I guess. 

%\addtocontents{toc}{\protect\vspace{0.5cm}}
%\appendix               % Resets chapter numbering to A, B, C... for appendices
%\raggedbottom\sloppy    % UWash thesis template has this
%\include{data_tables}
%\include{systematics}

%\include{bib}          % *.tex file for the Bibliography, or instead, use Latex4ZEUS!
\fixspacing
%\bibliography{bibliography,l4z_zeus,l4z_articles,l4z_conferences,l4z_misc,l4z_preprints,l4z_h1,l4z_temporary,syn}  % .bib files for the BibTex references
\bibliography{citations}
%\bibliographystyle{l4z_default}
\bibliographystyle{ieeetr}   % or below
%\bibliographystyle{unsrt}   % for using BibTex with references ordered as they appear

% See the thesis.cls file for the steps to convert AASTeX papers
% into thesis chapters


%%%%% BEGIN APPENDICES %%%%%%%%%%%%%%%%%%%%%%%%%%%%%%%%%%%%%%%%%%%%%%%%%%%%%%%%

%%%%%%%%%%%%%%%%%%%%%%%%%%%%%%%%%%%%%%%%%%%%%%%%%%%%%%%%%%%%%%%%%%%%%%%%%%%%%%%

\end{document}          % END THE DOCUMENT TEXT
